\documentclass[12pt,a4paper]{ctexart}
\usepackage{amsmath,amsthm,amssymb,bm,graphicx,hyperref,mathrsfs,geometry,float}

\title{XXX model in BS coordinate}
\author{Zhuan Ning}
\date{\today}
\linespread{1.25}

\geometry{left=1.25in,right=1.25in,top=1in,bottom=1in}

\numberwithin{equation}{section}

\hypersetup{colorlinks=true,linkcolor=blue,citecolor=blue}

\begin{document}

\maketitle

\section{Model}

The model used by Janik
\begin{equation}
    \mathcal{L}=R-\frac{1}{2}\nabla_\mu\phi\nabla^\mu\phi-V(\phi)
\end{equation}

where
\begin{equation}
    V(\phi)=-6\cosh(\frac{\phi}{\sqrt{3}})+b_4\phi^4
\end{equation}

with $b_4=-\frac{1}{5}$.

The model used by Tian
\begin{equation}
    \mathcal{L}=\frac{1}{2}(R-\nabla_\mu\psi\nabla^\mu\psi-2\tilde{V}(\psi))\sim R-\nabla_\mu\psi\nabla^\mu\psi-2\tilde{V}(\psi)
\end{equation}

so
\begin{equation}
    \psi=\frac{\phi}{\sqrt{2}},\quad\tilde{V}=\frac{V}{2}
\end{equation}

and
\begin{equation}
    \tilde{V}(\psi)=\frac{V(\phi)}{2}=\frac{V(\sqrt{2}\psi)}{2}=-3\cosh(\frac{\sqrt{6}\psi}{3})-\frac{2\psi^4}{5}
\end{equation}

\section{Spherical static solution}

the BS coordinate
\begin{equation}
    ds^2=\frac{L^2}{z^2}[-f(z)e^{-\chi(z)}dv^2-2e^{-\chi}dvdz+d\theta^2+\sin^2\theta d\varphi^2]
\end{equation}

\section{Qusinormal modes with azimuthal quantum number $l$}

the BS coordinate
\begin{equation}
    ds^2=\frac{L^2}{z^2}(-[f(v,z,\theta)e^{-\chi(v,z,\theta)}-e^{A(v,z,\theta)}\xi(v,z,\theta)^2]dv^2-2e^{-\chi}dvdz-2\xi e^Advd\theta+e^Ad\theta^2+e^{-A}\sin^2\theta d\varphi^2)
\end{equation}

one of the equations

\begin{equation}
    E_{z\theta}=z^2\left(\frac{e^{A+\chi} \xi'}{z^{2}}\right)' + (A+\chi)'_\theta + 2 \cot\theta A' + \frac{2 \chi_\theta}{z} - (A' A_\theta + \phi' \phi_\theta)
\end{equation}

linear perturbation on top of the spherical static background
\begin{equation}
    \xi=0,A=0,f_\theta=0,\chi_\theta=0,\phi_\theta=0,f_v=0,\chi_v=0,\phi_v=0
\end{equation}

\begin{equation}
    \delta E=z^2\left(\frac{e^{\chi} \delta\xi'}{z^{2}}\right)' + (\delta A+\delta\chi)'_\theta + 2 \cot\theta \delta A' + \frac{2 \delta\chi_\theta}{z} - \phi' \delta\phi_\theta
\end{equation}

\begin{equation}
    \begin{aligned}
        \Delta_2\delta\chi &= \frac{1}{\sin\theta}\partial_\theta(\sin\theta \delta\chi_\theta)\\
        \Delta_2\delta\phi &= \frac{1}{\sin\theta}\partial_\theta(\sin\theta \delta\phi_\theta)\\
        \delta\Xi &= \frac{1}{\sin\theta}\partial_\theta(\sin\theta \delta\xi)\\
        &= \delta\xi_\theta + \cot\theta\delta\xi\\
        \delta{a} &= \frac{1}{\sin\theta}\partial_\theta(\sin\theta (\delta A_\theta+2\cot\theta \delta A))\\
        &= \delta{A}_{\theta\theta} + 3 \cot\theta \delta{A}_\theta - 2 \delta{A}
    \end{aligned}
\end{equation}

\begin{equation}
    \frac{1}{\sin\theta}\partial_\theta(\sin\theta \delta E)=z^2\left(\frac{e^{\chi} \delta\Xi'}{z^{2}}\right)' + \delta a'+\Delta_2\delta\chi' + \frac{2 \Delta_2\delta\chi}{z} - \phi' \Delta_2\delta\phi
\end{equation}

\begin{equation}
    \delta\Psi=(\delta f,\delta\chi,\delta\Xi,\delta a,\delta\phi)\sim e^{-i\omega v}P_l(\cos\theta)\tilde{\Psi}(z)
\end{equation}

\begin{equation}
    \begin{aligned}
        \partial_v&\rightarrow-i\omega\\
        \partial_z&\rightarrow D\\
        \Delta_2&\rightarrow -l(l+1)
    \end{aligned}
\end{equation}

\subsection{boundary condition}

用演化的方法来求QNM其实就是在线性的层面做演化，所以要施加和时间演化一样的边界条件。

但是我们可以通过使用演化方程而不是约束方程（在这里是对于$f$和$\xi$），来避开施加类似于能量守恒和动量守恒这样的跟时间导数有关的边界条件。因此选择使用哪些方程来求QNM很关键，最简单的选择就是尽可能用演化方程。

我在一开始测试的时候，对于Schwarzschild-AdS黑洞，$f$和$\xi$全用约束方程，然后只施加source不变的边界条件，就能刚好得到Schwarzschild-AdS黑洞在纯标量扰动下的QNM。而后来测试把$f$和$\xi$换成演化方程就不行了，猜测是在之前的情况下，因为背景标量场为0，用约束方程的话可能刚好不起作用，就跟只是用标量场方程来求QNM得到了一样的结果。

\section{Time evolution}

\subsection{homogeneous case}

The dimension of spacetime is $d=1+3$, let $\kappa_4^2=8\pi G_4=1$. The Lagrangian density
\begin{equation}
    \mathcal{L}=\frac{1}{2}(R-\nabla_\mu\phi\nabla^\mu\phi-2V(\phi))
\end{equation}

the Einstein equation
\begin{equation}
    G_{\mu\nu}=R_{\mu\nu}-\frac{1}{2}Rg_{\mu\nu}=T_{\mu\nu}=T_{\mu\nu}^\phi
\end{equation}
\begin{equation}
    T_{\mu\nu}^\phi=\partial_\mu\phi\partial_\nu\phi-[\frac{1}{2}\partial_\rho\phi\partial^\rho\phi+V(\phi)]g_{\mu\nu}
\end{equation}

the equation of motion of scalar field
\begin{equation}
    \begin{aligned}
        \nabla_\mu\nabla^\mu\phi&=\frac{dV(\phi)}{d\phi}\\
        \frac{1}{\sqrt{-g}}\partial_\mu(\sqrt{-g}g^{\mu\nu}\partial_\nu\phi)&=V_\phi
    \end{aligned}
\end{equation}

the BS coordinate
\begin{equation}
    ds^2=\frac{L^2}{z^2}[-f(v,z)e^{-\chi(v,z)}dv^2-2e^{-\chi(v,z)}dvdz+d\Omega_2^2]
\end{equation}
\begin{equation}
    \sqrt{-g}=\frac{L^4}{z^4}e^{-\chi}\sin\theta
\end{equation}
\begin{equation}
    [g^{\mu\nu}]=\frac{z^2}{L^2}
        \begin{pmatrix}
            0&-e^\chi&&\\
            -e^\chi&fe^\chi&&\\
            &&1&&\\
            &&&\sin^{-2}\theta&\\
        \end{pmatrix}
\end{equation}
\begin{align}
    G_{vv}&=\frac{z f \left(f'-\dot{\chi}+z e^{-\chi }\right)-z \dot{f}+f^2 \left(z \chi'-3\right)}{z^2} \\
    G_{vz}&=\frac{z \left(f'+z e^{-\chi }\right)+f \left(z \chi'-3\right)}{z^2} \\
    G_{zz}&=\frac{2 \chi'}{z} \\
    G_{\theta\theta}&=-\frac{e^{\chi } \left(z \left(f' \left(z \chi'+4\right)-z \left(f''+2 \dot{\chi}'\right)\right)+f \left(z^2 \chi''-6\right)\right)}{2 z^2} \\
    G_{\varphi\varphi}&=\sin^2\theta G_{\theta\theta}
\end{align}
\begin{align}
    T_{vv}^\phi&=f \left(\frac{L^2 e^{-\chi } V(\phi )}{z^2}-\phi' \dot{\phi}\right)+\frac{1}{2} f^2 \phi'^2+\dot{\phi}^2 \\
    T_{vz}^\phi&=\frac{1}{2} f \phi'^2+\frac{L^2 e^{-\chi } V(\phi )}{z^2} \\
    T_{zz}^\phi&=\phi'^2 \\
    T_{\theta\theta}^\phi&=-\frac{1}{2} e^{\chi } \phi' \left(f \phi'-2 \dot{\phi}\right)-\frac{L^2 V(\phi )}{z^2} \\
    T_{\varphi\varphi}^\phi&=\sin^2\theta T_{\theta\theta}^\phi
\end{align}

the combinations of the Einstein equation
\begin{equation}
    G_{zz}=\frac{2\chi'}{z}=T_{zz}^\phi=\phi'^2=\Phi^2
\end{equation}
\begin{equation}
    \begin{aligned}
        G_{vv}-fG_{vz}&=\frac{-\dot{f}-f\dot{\chi}}{z}\\
        &=T_{vv}^\phi-fT_{vz}^\phi\\
        &=\dot{\phi}^2-f\dot{\phi}\phi'\\
        &=\Pi^2-f\Pi\Phi
    \end{aligned}
\end{equation}
\begin{equation}
    \begin{aligned}
        G_{vz}-\frac{f}{2}G_{zz}&=e^{-\chi}+\frac{f'}{z}-\frac{3f}{z^2}\\
        &=T_{vz}^\phi-\frac{f}{2}T_{zz}^\phi\\
        &=\frac{L^2 e^{-\chi } V(\phi )}{z^2}
    \end{aligned}
\end{equation}

where
\begin{equation}
    \begin{aligned}
        \Phi&=\phi'=\partial_z\phi\\
        \Pi&=\dot{\phi}=\partial_v\phi
    \end{aligned}
\end{equation}

the equation of motion of scalar field
\begin{equation}
    \begin{aligned}
        \partial_\mu(\sqrt{-g}g^{\mu\nu}\partial_\nu\phi)&=\partial_v(\sqrt{-g}g^{vz}\partial_z\phi)+\partial_z(\sqrt{-g}(g^{zv}\partial_v\phi+g^{zz}\partial_z\phi))\\
        &=\sqrt{-g}V_\phi\\
        \frac{\dot{\phi}}{z}-\dot{\phi}'+\frac{z^2}{2}\left(\frac{f\phi'}{z^2}\right)'&=\frac{L^2}{2z^2}e^{-\chi}V_\phi\\
        \Pi'-\frac{\Pi}{z}&=\frac{z^2}{2}\left(\frac{f\Phi}{z^2}\right)'-\frac{L^2}{2z^2}e^{-\chi}V_\phi
    \end{aligned}
\end{equation}

有如下两种演化方案：

第一种，$f$用演化方程：$\Phi$和$f$给初值，，从约束方程得到$\chi$和$\Pi$的初值，然后演化$\Phi$（即求它的时间导数），再演化$f$（注意，
令$f=\tilde{f}e^{-\chi}$，可以把$f$的演化方程和$\chi$退耦）

第二种，$f$用约束方程：$\Phi$给初值，从约束方程得到$\chi$,$f$和$\Pi$的初值，边界条件分别为$\chi|_B=0$，和从$f$的演化方程渐近展开得到的一个等式（即对$f$的能量约束），然后就能演化$\Phi$了

在Mathematica程序里，$\delta$是$\chi$，A是$\Phi$，P是$\Pi$，diz是对z积分，dz是对z求导。

\subsection{homogeneous case with superradiation}

\begin{equation}
    I=\int d^{4}y\sqrt{-g}[\frac{1}{2}R-\frac{1}{2}(\partial\phi)^{2}-V(\phi)-\frac{1}{4}F(\phi)F^{\sigma\rho}F_{\sigma\rho}]
\end{equation}

\begin{equation}
    e^{-\chi}+\frac{f^{\prime}}{z}-\frac{3f}{z^{2}}=\frac{L^{2}}{z^{2}}e^{-\chi}V+\frac{1}{2}\frac{z^{2}}{L^{2}}e^{\chi}FA_{v}^{\prime2}
\end{equation}

\begin{equation}
    \frac{1}{z}[-\dot{f}-f\dot{\chi}]=\Pi^{2}-f\Pi\Phi
\end{equation}

\begin{equation}
    \dot{f}|_{r_h}=-\frac{1}{r_h}\dot{\phi}^2<0
\end{equation}

\begin{equation}
    f'|_{r_h}=-\frac{1}{r_h}e^{-\chi}+r_h e^{-\chi}V(\phi)+\frac{1}{2r_h^3}e^{\chi}F(\phi)A^{\prime2}
\end{equation}

\begin{equation}
    V(\phi)=-6+\frac{1}{2}m^2\phi^2<0
\end{equation}

\begin{equation}
    \dot{z}_{h}=-\frac{\dot{f}_{h}}{f_{h}^{\prime}}
\end{equation}

\begin{equation}
    \dot{z}_{h}=-\frac{\dot{r}_h}{r_h^2}
\end{equation}

\begin{equation}
    \partial_t(r_h)=\frac{2\dot{\phi}^2}{2r_h^2e^{-\chi}-2r_h^4 e^{-\chi}(2\Lambda+m^2\phi^2)-e^{\chi}F(\phi)A^{\prime2}}
\end{equation}

\subsection{inhomogeneous case with planar topology}

\subsubsection{equation}

The dimension of spacetime is $d=1+3$, let $\kappa_4^2=8\pi G_4=1$. The Lagrangian density
\begin{equation}
    \mathcal{L}=R-\frac{1}{2}\nabla_\mu\phi\nabla^\mu\phi-V(\phi)
\end{equation}

the Einstein equation
\begin{equation}
    G_{\mu\nu}=R_{\mu\nu}-\frac{1}{2}Rg_{\mu\nu}=T_{\mu\nu}=T_{\mu\nu}^\phi
\end{equation}
\begin{equation}
    T_{\mu\nu}^\phi=\frac{1}{2}(\partial_\mu\phi\partial_\nu\phi-[\frac{1}{2}\partial_\rho\phi\partial^\rho\phi+V(\phi)]g_{\mu\nu})
\end{equation}
\begin{equation}
    E_{\mu\nu}=G_{\mu\nu}-T_{\mu\nu}
\end{equation}

the equation of motion of scalar field
\begin{equation}
    \begin{aligned}
        \nabla_\mu\nabla^\mu\phi&=\frac{dV(\phi)}{d\phi}\\
        \frac{1}{\sqrt{-g}}\partial_\mu(\sqrt{-g}g^{\mu\nu}\partial_\nu\phi)&=V_\phi
    \end{aligned}
\end{equation}

the BS coordinate
\begin{equation}
    ds^2=\frac{L^2}{z^2}(-[f(v,z,x)e^{-\chi(v,z,x)}-e^{A(v,z,x)}\xi(v,z,x)^2]dv^2-2e^{-\chi}dvdz-2\xi e^Advdx+e^Adx^2+e^{-A}dy^2)
\end{equation}

$E_{zz}$ gives the constraint equation of $\chi$:
\begin{equation}
    \chi'=\frac{z}{4}(A'^2+\phi'^2)
\end{equation}

\begin{equation}
    z\left(\frac{\chi}{z^2}\right)'+2\frac{\chi}{z^2}=\frac{1}{4}(A'^2+\phi'^2)
\end{equation}

$E_{zx}$ gives the constraint equation of $\xi$:
\begin{equation}
    \begin{aligned}
        \left(\frac{e^{A+\chi}}{z^2}\xi'\right)'&=-\frac{1}{z^{2}}[(A+\chi)^{\prime}_{x}+\frac{2\chi_{x}}{z}-(A^{\prime}A_{x}+\phi'\phi_{x})]\\
        \xi'&=\left(\frac{e^{A+\chi}}{z^2}\xi'\right)\frac{z^2}{e^{A+\chi}}
    \end{aligned}
\end{equation}

$E_{vz}-\frac{f}{2}E_{zz}+\xi E_{zx}$ gives the constraint equation of $f$:
\begin{equation}
    \begin{aligned}
        \mathrm{COM}&=\frac{e^{A+\chi}}{4}\xi^{\prime2}+\frac{\xi_{x}}{z}-\frac{\xi_{x}^{\prime}}{2}-\frac{e^{-A-\chi}}{4}(2A_x\chi_x+\chi_{x}^2-\phi_x^2)+\frac{(e^{-A-\chi})_{xx}}{2}\\
        &=\frac{e^{A+\chi}}{4}\xi^{\prime2}+\frac{\xi_{x}}{z}-\frac{\xi_{x}^{\prime}}{2}+\frac{e^{-A-\chi}}{2}[A_x^2+A_x\chi_x+\frac{1}{2}\chi_{x}^2+\frac{1}{2}\phi_x^2-(A_{xx}+\chi_{xx})]
    \end{aligned}
\end{equation}
\begin{equation}
    \begin{aligned}
        (\frac{f}{z^{3}})^{\prime}&=\frac{\xi_{x}}{z^{3}}+\frac{L^{2}}{2z^{4}}e^{-\chi}V+\frac{\mathrm{COM}}{z^{2}}\\
        &=(\frac{1+z^2\tilde{f}}{z^{3}})^{\prime}\\
        z(\frac{\tilde{f}}{z})^{\prime}&=\tilde{f}'-\frac{\tilde
        {f}}{z}\\
        &=\frac{\xi_{x}}{z^2}+\frac{1}{z^3}(3+\frac{L^{2}}{2}e^{-\chi}V)+\frac{\mathrm{COM}}{z}
    \end{aligned}
\end{equation}

\begin{equation}
    \begin{aligned}
        (\frac{f}{z^{3}})^{\prime}&=\frac{\xi_{x}}{z^{3}}+\frac{L^{2}}{2z^{4}}e^{-\chi}V+\frac{\mathrm{COM}}{z^{2}}\\
        zf'-3f&=z\xi_{x}+\frac{L^{2}}{2}e^{-\chi}V+z^2\mathrm{COM}\\
        &=2z\xi_{x}+\frac{L^{2}}{2}e^{-\chi}V+z^2(\mathrm{COM}-\frac{\xi_x}{z})\\
        &=2z\xi_{x}+\frac{L^{2}}{2}e^{-\chi}V+z^2\mathrm{COM}_f\\
        \mathrm{COM}_f&=\mathrm{COM}-\frac{\xi_x}{z}\\
        &=\frac{e^{A+\chi}}{4}\xi^{\prime2}-\frac{\xi_{x}^{\prime}}{2}-\frac{e^{-A-\chi}}{4}(2A_x\chi_x+\chi_{x}^2-\phi_x^2)+\frac{(e^{-A-\chi})_{xx}}{2}\\
        &=\frac{e^{A+\chi}}{4}\xi^{\prime2}-\frac{\xi_{x}^{\prime}}{2}+\frac{e^{-A-\chi}}{2}[A_x^2+A_x\chi_x+\frac{1}{2}\chi_{x}^2+\frac{1}{2}\phi_x^2-(A_{xx}+\chi_{xx})]
    \end{aligned}
\end{equation}

$e^{-A-\chi}E_{xx}-e^{A-\chi}E_{yy}$ gives the evolution equation of $A$:
\begin{equation}
    \begin{aligned}
        z(\frac{\dot{A}}{z})^{\prime}&=\dot{A}^{\prime}-\frac{\dot{A}}{z}\\
        &=\frac{z^{2}}{2}(\frac{fA^{\prime}-\xi A_{x}}{z^{2}})^{\prime}+\frac{(e^{-A-\chi}A_{x}-\xi A^{\prime})_{x}}{2}+\mathrm{COM}\\
        &=\frac{(fA^{\prime}-\xi A_{x})'}{2}-\frac{fA^{\prime}-\xi A_{x}}{z}+\frac{(e^{-A-\chi}A_{x}-\xi A^{\prime})_{x}}{2}+\mathrm{COM}\\
        &=\frac{(fA^{\prime}-\xi A_{x})'}{2}-\frac{fA^{\prime}-\xi A_{x}}{z}-\frac{(\xi A^{\prime})_{x}}{2}\\
        &\quad+\frac{e^{A+\chi}}{4}\xi^{\prime2}+\frac{\xi_{x}}{z}-\frac{\xi_{x}^{\prime}}{2}+\frac{e^{-A-\chi}}{2}[\frac{1}{2}\chi_{x}^2+\frac{1}{2}\phi_x^2-\chi_{xx}]
    \end{aligned}
\end{equation}

\begin{equation}
    z(\frac{\dot{A}}{z})^{\prime}=\frac{(fA^{\prime}-\xi A_{x})'}{2}-\frac{fA^{\prime}-\xi A_{x}}{z}+\frac{(e^{-A-\chi}A_{x}-\xi A^{\prime})_{x}}{2}+\mathrm{COM}
\end{equation}

\begin{equation}
    \dot{A}^{\prime}-\frac{\dot{A}}{z}=\frac{(fA^{\prime}-\xi A_{x})'}{2}-\frac{fA^{\prime}-\xi A_{x}}{z}+\frac{(e^{-A-\chi}A_{x}-\xi A^{\prime})_{x}}{2}+\mathrm{COM}
\end{equation}

the scalar equation gives the evolution equation of $\phi$:
\begin{equation}
    \begin{aligned}
        z(\frac{\dot{\phi}}{z})^{\prime}&=\dot{\phi}^{\prime}-\frac{\dot{\phi}}{z}\\
        &=\frac{z^{2}}{2}(\frac{f\phi'-\xi\phi_{x}}{z^{2}})^{\prime}+\frac{(e^{-A-\chi}\phi_{x}-\xi\phi')_{x}}{2}-\frac{L^{2}}{2z^{2}}e^{-\chi}V_\phi\\
        &=\frac{(f\phi'-\xi\phi_{x})'}{2}-\frac{f\phi'-\xi\phi_{x}}{z}+\frac{(e^{-A-\chi}\phi_{x}-\xi\phi')_{x}}{2}-\frac{L^{2}}{2z^{2}}e^{-\chi}V_\phi
    \end{aligned}
\end{equation}

\begin{equation}
    z(\frac{\dot{\phi}}{z})^{\prime}=\frac{(f\phi'-\xi\phi_{x})'}{2}-\frac{f\phi'-\xi\phi_{x}}{z}+\frac{(e^{-A-\chi}\phi_{x}-\xi\phi')_{x}}{2}-\frac{L^{2}}{2z^{2}}e^{-\chi}V_\phi
\end{equation}

\begin{equation}
    \dot{\phi}^{\prime}-\frac{\dot{\phi}}{z}=\frac{(f\phi'-\xi\phi_{x})'}{2}-\frac{f\phi'-\xi\phi_{x}}{z}+\frac{(e^{-A-\chi}\phi_{x}-\xi\phi')_{x}}{2}-\frac{L^{2}}{2z^{2}}e^{-\chi}V_\phi
\end{equation}

$E_{xn}=E_{xv}-fE_{xz}+\xi E_{xx}$ gives the evolution equation of $\xi$:
\begin{equation}
    \begin{aligned}
        0&=\frac{2}{z}(f_{x}+f\chi_{x})+(2\xi e^{A+\chi}\xi^{\prime}-f^{\prime}+f[A^{\prime}+\chi^{\prime}])_{x}-\xi(e^{A+\chi}\xi^{\prime})_{x}-\xi(A+\chi)_{xx}\\
        &\quad+A_{x}d_{n}A+\phi_{x}d_n\phi-(\dot{A}+\dot{\chi})_{x}+(e^{A+\chi}\xi^{\prime})_{v}\\
        &=\frac{2}{z}(f_{x}+f\chi_{x})+2\xi_{x}e^{A+\chi}\xi^{\prime}+(f[A^{\prime}+\chi^{\prime}]-f^{\prime})_{x}+\xi(e^{A+\chi}\xi^{\prime})_{x}-\xi(A+\chi)_{xx}\\
        &\quad+A_{x}d_{n}A+\phi_{x}d_n\phi-(\dot{A}+\dot{\chi})_{x}+(e^{A+\chi}\xi^{\prime})_{v}\\
        &=\frac{2}{z}(f_{x}+f\chi_{x})+(2\xi e^{A+\chi}\xi^{\prime}-f^{\prime}+f[A^{\prime}+\chi^{\prime}]-[\dot{A}+\dot{\chi}])_{x}-\xi(e^{A+\chi}\xi^{\prime}+A_x+\chi_x)_{x}\\
        &\quad+A_{x}d_{n}A+\phi_{x}d_n\phi+(e^{A+\chi}\xi^{\prime})_{v}
    \end{aligned}
\end{equation}

where $d_{n}:=\partial_{v}-f\partial_{z}+\xi\partial_{x}$.

$E_{vn}=E_{vv}-fE_{vz}+\xi E_{vx}$ gives the evolution equation of $f$:
\begin{equation}
    \begin{aligned}
        0&=(e^{-A-\chi}f_{x}-[f\xi]^{\prime}+\xi^{2}e^{A+\chi}\xi^{\prime}-\xi[d_{n}A+d_{n}\chi]-2d_{n}\xi)_{x}-\xi_x(\dot{A}+\dot{\chi})\\
        &\quad-\frac{2}{z}(\dot{f}+f\dot{\chi})-\dot{A}d_{n}A-\dot{\phi}d_n\phi+\xi(e^{A+\chi}\xi^{\prime})_{v}
    \end{aligned}
\end{equation}

$\chi$, $f$和$\xi$的时间导数的计算：

$\partial_v E_{zz}$ gives the evolution equation of $\chi$:
\begin{equation}
    \dot{\chi}'=\frac{z}{2}(A'\dot{A}'+\phi'\dot{\phi}')
\end{equation}

$E_{xn}=E_{xv}-fE_{xz}+\xi E_{xx}$ gives the evolution equation of $\xi$:
\begin{equation}
    \begin{aligned} 
        (e^{A+\chi}\xi^{\prime})_{v}&=-\frac{2}{z}(f_{x}+f\chi_{x})-(2\xi e^{A+\chi}\xi^{\prime}-f^{\prime}+f[A^{\prime}+\chi^{\prime}]-[\dot{A}+\dot{\chi}])_{x}+\xi(e^{A+\chi}\xi^{\prime}+A_{x}+\chi_{x})_{x}\\
        &\quad-A_{x}d_{n}A-\phi_{x}d_n\phi\\
        (e^{A+\chi}\xi^{\prime})_{v}&=e^{A+\chi}[\xi^{\prime}(\dot{A}+\dot{\chi})+\dot{\xi}']\\
        \dot{\xi}'&=e^{-A-\chi}(e^{A+\chi}\xi^{\prime})_{v}-\xi^{\prime}(\dot{A}+\dot{\chi})
    \end{aligned}
\end{equation}

$E_{vn}=E_{vv}-fE_{vz}+\xi E_{vx}$ gives the evolution equation of $f$:
\begin{equation}
    \begin{aligned}
        \dot{f}&=\frac{z}{2}[(e^{-A-\chi}f_{x}-[f\xi]^{\prime}+\xi^{2}e^{A+\chi}\xi^{\prime}-\xi[d_{n}A+d_{n}\chi]-2d_{n}\xi)_{x}-\xi_x(\dot{A}+\dot{\chi})\\
        &\quad-\dot{A}d_{n}A-\dot{\phi}d_n\phi+\xi(e^{A+\chi}\xi^{\prime})_{v}]-f\dot{\chi}
    \end{aligned}
\end{equation}

alternatively, $\partial_v E_{zx}$ and $\partial_v (E_{vz}-\frac{f}{2}E_{zz}+\xi E_{zx})$ gives
\begin{equation}
    \begin{aligned}
        \left(\frac{(e^{A+\chi}\xi')_v}{z^2}\right)'&=-\frac{1}{z^{2}}[(\dot{A}+\dot{\chi})^{\prime}_{x}+\frac{2\dot{\chi}_{x}}{z}
        -(\dot{A}_{x}A^{\prime}+A_{x}\dot{A}^{\prime}+\dot{\phi}_{x}\phi'+\phi_{x}\dot{\phi}')]\\
        \dot{\xi}'&=z^2e^{-A-\chi}\frac{(e^{A+\chi}\xi^{\prime})_{v}}{z^2}-\xi^{\prime}(\dot{A}+\dot{\chi})
    \end{aligned}
\end{equation}

\begin{equation}
    \begin{aligned}
        \partial_v\mathrm{COM}&=\frac{e^{A+\chi}(\dot{A}+\dot{\chi})}{4}\xi^{\prime2}+\frac{e^{A+\chi}}{2}\xi^{\prime}\dot{\xi}'+\frac{\dot{\xi}_{x}}{z}-\frac{\dot{\xi}_{x}^{\prime}}{2}\\
        &\quad+\frac{e^{-A-\chi}(\dot{A}+\dot{\chi})}{4}(2A_x\chi_x+\chi_{x}^2-\phi_x^2)-\frac{e^{-A-\chi}}{4}(2\dot{A}_x\chi_x+2A_x\dot{\chi}_x+2\chi_{x}\dot{\chi}_{x}-2\phi_x\dot{\phi}_x)\\
        &\quad-\frac{(e^{-A-\chi}(\dot{A}+\dot{\chi}))_{xx}}{2}\\
        &=\frac{e^{A+\chi}}{4}\xi'[\xi'(\dot{A}+\dot{\chi})+2\dot{\xi}']+\frac{\dot{\xi}_{x}}{z}-\frac{\dot{\xi}_{x}^{\prime}}{2}\\
        &\quad-\frac{e^{-A-\chi}}{4}[-(\dot{A}+\dot{\chi})(2A_x\chi_x+\chi_{x}^2-\phi_x^2)+(2\dot{A}_x\chi_x+2A_x\dot{\chi}_x+2\chi_{x}\dot{\chi}_{x}-2\phi_x\dot{\phi}_x)]\\
        &\quad-\frac{(e^{-A-\chi}(\dot{A}+\dot{\chi}))_{xx}}{2}
    \end{aligned}
\end{equation}
\begin{equation}
    \begin{aligned}
        z(\frac{\dot{\tilde{f}}}{z})^{\prime}&=\dot{\tilde{f}}'-\frac{\dot{\tilde
        {f}}}{z}\\
        &=(-\frac{L^{2}}{2}e^{-\chi}\dot{\chi}V+\frac{L^{2}}{2}e^{-\chi}V_\phi\dot{\phi})\frac{1}{z^3}+\frac{\dot{\xi}_{x}}{z^2}+\frac{\partial_v\mathrm{COM}}{z}
    \end{aligned}
\end{equation}

\subsubsection{asymptotic behavior}

Assumning $L=1$, we expand the fields near the boundary ($z=0$):
\begin{equation}
    \begin{aligned}
        f&=f_0+f_1 z+f_2 z^2+f_3 z^3+O(z^4)\\
        \chi&=\chi_0+\chi_1 z+\chi_2 z^2+\chi_3 z^3+O(z^4)\\
        \xi&=\xi_0+\xi_1 z+\xi_2 z^2+\xi_3 z^3+O(z^4)\\
        A&=A_0+A_1 z+A_2 z^2+A_3 z^3+O(z^4)\\
        \phi&=\phi_0+\phi_1 z+\phi_2 z^2+\phi_3 z^3+O(z^4)\\
    \end{aligned}
\end{equation}

both $f_i$ are function of $(v,x)$, substituting them into equations of motion, we can get
\begin{equation}
    \chi_1=0,\phi_0=0,\xi_1=e^{-A_0-\chi_0}\partial_x\chi_0,f_0=e^{-\chi_0},f_1=-\partial_x\xi_0
\end{equation}

after letting $\chi_0=0,\xi_0=0$, we can get
\begin{equation}
    \xi_1=0,f_0=1,f_1=0
\end{equation}
\begin{equation}
    A_1=\partial_v A_0,A_2=0
\end{equation}

after letting $A_0=0$, we can get
\begin{equation}
    A_1=0
\end{equation}
\begin{equation}
    \xi_2=0,f_2=\chi_2,\chi_2=\frac{\phi_1^2}{8},\chi_3=\frac{\phi_1\phi_2}{3}
\end{equation}
\begin{equation}
    \phi_3=\partial_v\phi_2-\frac{119\phi_1^3}{360}-\frac{\partial_x^2\phi_1}{2}
\end{equation}
\begin{equation}
    \partial_v\xi_3=\frac{\partial_x f_3}{3}-\partial_x A_3-\frac{2\phi_1\partial_x\phi_2}{9}+\frac{\phi_2\partial_x\phi_1}{9}-\frac{\partial_v\phi_1\partial_x\phi_1}{4}+\frac{\phi_1\partial_{vx}\phi_1}{12}
\end{equation}
\begin{equation}
    \partial_v f_3=\frac{3\partial_x\xi_3}{2}+\frac{\phi_1\partial_v\phi_2}{6}+\frac{2\phi_2\partial_v\phi_1}{3}-\frac{(\partial_v\phi_1)^2}{2}+\frac{(\partial_x\phi_1)^2}{8}+\frac{\phi_1\partial_x^2\phi_1}{8}
\end{equation}

after letting $\phi_1(v,x)=\phi_1$, we can get
\begin{equation}
    \begin{aligned}
        f&=1+\frac{\phi_1^2}{8}z^2+f_3 z^3+O(z^4)\\
        \chi&=\frac{\phi_1^2}{8}z^2+\frac{\phi_1\phi_2}{3}z^3+O(z^4)\\
        \xi&=\xi_3 z^3+O(z^4)\\
        A&=A_3 z^3+O(z^4)\\
        \phi&=\phi_1 z+\phi_2 z^2+(\partial_v\phi_2-\frac{119\phi_1^3}{360})z^3+O(z^4)\\
    \end{aligned}
\end{equation}
\begin{equation}
    \partial_v\xi_3=\frac{\partial_x f_3}{3}-\partial_x A_3-\frac{2\phi_1\partial_x\phi_2}{9}=\partial_x(\frac{ f_3}{3}-A_3-\frac{2\phi_1\phi_2}{9})
\end{equation}
\begin{equation}
    \partial_v f_3=\frac{3\partial_x\xi_3}{2}+\frac{\phi_1\partial_v\phi_2}{6}
\end{equation}

\subsubsection{subtraction}

let
\begin{equation}
    \begin{aligned}
        f&=1+\frac{\phi_1^2}{8}z^2+z^3\tilde{f}\\
        \chi&=\frac{\phi_1^2}{8}z^2+z^3\tilde{\chi}\\
        \xi&=z^3\tilde{\xi}\\
        A&=z^3\tilde{A}\\
        \phi&=\phi_1 z+z^2\tilde{\phi}
    \end{aligned}
\end{equation}

the equations become
\begin{equation}
    \begin{aligned}
        z(z^2\tilde{\chi}'+3z\tilde{\chi}+\frac{\phi_1^2}{4})&=\frac{z}{4}((z^3\tilde{A}'+3z^2\tilde{A})^2+(z^2\tilde{\phi}'+2z\tilde{\phi}+\phi_1)^2)\\
        z^2\tilde{\chi}'+3z\tilde{\chi}+\frac{\phi_1^2}{4}&=\frac{1}{4}((z^3\tilde{A}'+3z^2\tilde{A})^2+(z^2\tilde{\phi}'+2z\tilde{\phi})^2+2(z^2\tilde{\phi}'+2z\tilde{\phi})\phi_1)+\frac{\phi_1^2}{4}\\
        z\tilde{\chi}'+3\tilde{\chi}&=\frac{1}{4}(z^3(z\tilde{A}'+3\tilde{A})^2+z(z\tilde{\phi}'+2\tilde{\phi})^2+2(z\tilde{\phi}'+2\tilde{\phi})\phi_1)
    \end{aligned}
\end{equation}
\begin{equation}
    \begin{aligned}
        A'&=z^2(z\tilde{A}'+3\tilde{A})\\
        \phi'&=z(z\tilde{\phi}'+2\tilde{\phi})+\phi_1
    \end{aligned}
\end{equation}
\begin{equation}
    \begin{aligned}
        \left(e^{A+\chi}(z\tilde{\xi}'+3\tilde{\xi})\right)'&=-[z(\tilde{A}+\tilde{\chi})^{\prime}_{x}+3(\tilde{A}+\tilde{\chi})_{x}+2\tilde{\chi}_{x}-(z\tilde{A}_{x}A'+\tilde{\phi}_{x}\phi')]\\
        &=-(z\tilde{A}'+3\tilde{A}+z\tilde{\chi}'+3\tilde{\chi})_{x}-2\tilde{\chi}_{x}+(z^3\tilde{A}_{x}(z\tilde{A}'+3\tilde{A})+\tilde{\phi}_{x}\phi')\\
        z\tilde{\xi}'+3\tilde{\xi}&=\left(e^{A+\chi}(z\tilde{\xi}'+3\tilde{\xi})\right)e^{-A-\chi}
    \end{aligned}
\end{equation}
\begin{equation}
    \begin{aligned}
        \mathrm{COM}&=z^4\frac{e^{A+\chi}}{4}(z\tilde{\xi}'+3\tilde{\xi})^2+z^2\tilde{\xi}_{x}-\frac{z^2(z\tilde{\xi}'+3\tilde{\xi})_{x}}{2}-z^4\frac{e^{-A-\chi}}{4}(2z^2\tilde{A}_x\tilde{\chi}_x+z^2\tilde{\chi}_{x}^2-\tilde{\phi}_x^2)\\
        &\quad-z^3\frac{(e^{-A-\chi}(\tilde{A}+\tilde{\chi})_x)_{x}}{2}
    \end{aligned}
\end{equation}
\begin{equation}
    \begin{aligned}
        \mathrm{COM}_s&=\frac{\mathrm{COM}}{z^2}\\
        &=z^2\frac{e^{A+\chi}}{4}(z\tilde{\xi}'+3\tilde{\xi})^2+\tilde{\xi}_{x}-\frac{(z\tilde{\xi}'+3\tilde{\xi})_{x}}{2}-z^2\frac{e^{-A-\chi}}{4}(2z^2\tilde{A}_x\tilde{\chi}_x+z^2\tilde{\chi}_{x}^2-\tilde{\phi}_x^2)\\
        &\quad-z\frac{(e^{-A-\chi}(\tilde{A}+\tilde{\chi})_x)_{x}}{2}
    \end{aligned}
\end{equation}
\begin{equation}
    \begin{aligned}
        (\frac{1}{z^{3}}+\frac{\phi_1^2}{8z}+\tilde{f})^{\prime}&=\frac{L^{2}}{2z^{4}}e^{-\chi}V+\frac{\xi_{x}}{z^{3}}+\frac{\mathrm{COM}}{z^{2}}\\
        \tilde{f}'&=(3+\frac{L^{2}}{2}e^{-\chi}V)\frac{1}{z^4}+\frac{\phi_1^2}{8z^2}+\tilde{\xi}_{x}+\mathrm{COM}_s
    \end{aligned}
\end{equation}
\begin{equation}
    \begin{aligned}
        z(z^2\dot{\tilde{A}})^{\prime}&=z^2(z\dot{\tilde{A}}^{\prime}+2\dot{\tilde{A}})\\
        &=\frac{z^{2}}{2}(\frac{fA^{\prime}-\xi A_{x}}{z^{2}})^{\prime}+\frac{(e^{-A-\chi}A_{x}-\xi A^{\prime})_{x}}{2}+\mathrm{COM}\\
        z\dot{\tilde{A}}^{\prime}+2\dot{\tilde{A}}&=\frac{1}{2}(f(z\tilde{A}'+3\tilde{A})-z^4\tilde{\xi}\tilde{A}_{x})^{\prime}+\frac{z(e^{-A-\chi}\tilde{A}_{x}-\tilde{\xi} A^{\prime})_{x}}{2}+\mathrm{COM}_s
    \end{aligned}
\end{equation}
\begin{equation}
    \begin{aligned}
        z(z\dot{\tilde{\phi}})^{\prime}&=z(z\dot{\tilde{\phi}}^{\prime}+\dot{\tilde{\phi}})\\
        &=\frac{z^{2}}{2}(\frac{f\phi'-\xi\phi_{x}}{z^{2}})^{\prime}+\frac{(e^{-A-\chi}\phi_{x}-\xi\phi')_{x}}{2}-\frac{L^{2}}{2z^{2}}e^{-\chi}V_\phi\\
        &=\frac{(f\phi'-\xi\phi_{x})'}{2}-\frac{f\phi'-\xi\phi_{x}}{z}+\frac{(e^{-A-\chi}\phi_{x}-\xi\phi')_{x}}{2}-\frac{L^{2}}{2z^{2}}e^{-\chi}V_\phi\\
        z\dot{\tilde{\phi}}^{\prime}+\dot{\tilde{\phi}}&=\frac{z}{2}(\tilde{\phi}'+(\frac{\phi_1^2}{8}+z\tilde{f})\phi'-z^3\tilde{\xi}\tilde{\phi}_{x})^{\prime}+\frac{z}{2}(\frac{2\tilde{\phi}}{z}+\frac{\phi_1}{z^2})'+z\frac{(e^{-A-\chi}\tilde{\phi}_{x}-z\tilde{\xi}\phi')_{x}}{2}-\frac{L^{2}}{2z^{3}}e^{-\chi}V_\phi\\
        &=\frac{z}{2}([\tilde{\phi}'+(\frac{\phi_1^2}{8}+z\tilde{f})\phi'-z^3\tilde{\xi}\tilde{\phi}_{x}]^{\prime}+(e^{-A-\chi}\tilde{\phi}_{x}-z\tilde{\xi}\phi')_{x})+\tilde{\phi}'-\frac{\tilde{\phi}}{z}-\frac{\phi_1}{z^2}-\frac{L^{2}}{2z^{3}}e^{-\chi}V_\phi
    \end{aligned}
\end{equation}
\begin{equation}
    z\dot{\tilde{\chi}}'+3\dot{\tilde{\chi}}=\frac{1}{2}(z^3(z\tilde{A}'+3\tilde{A})(z\dot{\tilde{A}}'+3\dot{\tilde{A}})+z(z\tilde{\phi}'+2\tilde{\phi})(z\dot{\tilde{\phi}}'+2\dot{\tilde{\phi}})+(z\dot{\tilde{\phi}}'+2\dot{\tilde{\phi}})\phi_1)
\end{equation}
\begin{equation}
    \begin{aligned}
        (e^{A+\chi}(z\tilde{\xi}'+3\tilde{\xi}))_{v}&=-2(\tilde{f}_{x}+f\tilde{\chi}_{x})\\
        &\quad-(2ze^{A+\chi}\tilde{\xi}\xi^{\prime}-(z\tilde{f}'+3\tilde{f})+f[z\tilde{A}'+3\tilde{A}+z\tilde{\chi}'+3\tilde{\chi}]+z^2\tilde{f}\frac{\phi_1^2}{4}-z[\dot{\tilde{A}}+\dot{\tilde{\chi}}])_{x}\\
        &\quad+z\tilde{\xi}(e^{A+\chi}\xi^{\prime}+A_{x}+\chi_{x})_{x}-z\tilde{A}_{x}d_{n}A-\tilde{\phi}_{x}d_n\phi\\
        (e^{A+\chi}(z\tilde{\xi}'+3\tilde{\xi}))_{v}&=e^{A+\chi}[z^3(z\tilde{\xi}'+3\tilde{\xi})(\dot{\tilde{A}}+\dot{\tilde{\chi}})+(z\dot{\tilde{\xi}}'+3\dot{\tilde{\xi}})]\\
        z\dot{\tilde{\xi}}'+3\dot{\tilde{\xi}}&=e^{-A-\chi}(e^{A+\chi}(z\tilde{\xi}'+3\tilde{\xi}))_{v}-z^3(z\tilde{\xi}'+3\tilde{\xi})(\dot{\tilde{A}}+\dot{\tilde{\chi}})
    \end{aligned}
\end{equation}
\begin{equation}
    \begin{aligned}
        \dot{\tilde{f}}&=\frac{z}{2}[(e^{-A-\chi}\tilde{f}_{x}-f'\tilde{\xi}+z^3e^{A+\chi}\tilde{\xi}^{2}\xi^{\prime}-\tilde{\xi}[d_{n}A+d_{n}\chi])_{x}-z^3\tilde{\xi}_x(\dot{\tilde{A}}+\dot{\tilde{\chi}})\\
        &\quad-\dot{\tilde{A}}d_{n}A+z^2\tilde{\xi}(e^{A+\chi}(z\tilde{\xi}'+3\tilde{\xi}))_{v}]-f\dot{\tilde{\chi}}-\frac{1}{2}([f(z\tilde{\xi}'+3\tilde{\xi})+2\frac{d_{n}\xi}{z^2}]_x+\dot{\tilde{\phi}}d_n\phi)
    \end{aligned}
\end{equation}

\begin{equation}
    \begin{aligned}
        (e^{A+\chi}(z\tilde{\xi}'+3\tilde{\xi}))_{v}'&=-(z\dot{\tilde{A}}'+3\dot{\tilde{A}}+z\dot{\tilde{\chi}}'+3\dot{\tilde{\chi}})_{x}-2\dot{\tilde{\chi}}_{x}\\
        &\quad+(z^3\dot{\tilde{A}}_{x}(z\tilde{A}'+3\tilde{A})+z^3\tilde{A}_{x}(z\dot{\tilde{A}}'+3\dot{\tilde{A}})+\dot{\tilde{\phi}}_{x}\phi'+z\tilde{\phi}_{x}(z\dot{\tilde{\phi}}'+2\dot{\tilde\phi}))\\
        z\dot{\tilde{\xi}}'+3\dot{\tilde{\xi}}&=e^{-A-\chi}(e^{A+\chi}(z\tilde{\xi}'+3\tilde{\xi}))_{v}-z^3(z\tilde{\xi}'+3\tilde{\xi})(\dot{\tilde{A}}+\dot{\tilde{\chi}})
    \end{aligned}
\end{equation}

transforamtion

\begin{equation}
    X:=e^{-\chi}=1-\frac{\phi_1^2}{8}z^2+z^3\tilde{X}
\end{equation}

\begin{equation}
    \frac{X'}{X}=-\frac{z}{4}(A'^2+\phi'^2)
\end{equation}

\begin{equation}
    \begin{aligned}
        X'&=-\frac{zX}{4}(A'^2+\phi'^2)\\
        z^2\tilde{X}'+3z\tilde{X}-\frac{\phi_1^2}{4}&=-\frac{1}{4}(A'^2+\phi'^2)(1-\frac{\phi_1^2}{8}z^2+z^3\tilde{X})\\
        z^2\tilde{X}'+3z\tilde{X}+\frac{z^3}{4}(A'^2+\phi'^2)\tilde{X}&=-\frac{1}{4}(A'^2+\phi'^2)(1-\frac{\phi_1^2}{8}z^2)+\frac{\phi_1^2}{4}\\
        z^2\tilde{X}'+3z\tilde{X}+\frac{z^3}{4}(A'^2+\phi'^2)\tilde{X}&=-\frac{1}{4}((z^3\tilde{A}'+3z^2\tilde{A})^2+(z^2\tilde{\phi}'+2z\tilde{\phi}+\phi_1)^2)(1-\frac{\phi_1^2}{8}z^2)+\frac{\phi_1^2}{4}\\
        z\tilde{X}'+3\tilde{X}+\frac{z^2}{4}(A'^2+\phi'^2)\tilde{X}&=-\frac{1}{4}(z^3(z\tilde{A}'+3\tilde{A})^2+z(z\tilde{\phi}'+2\tilde{\phi})^2+2(z\tilde{\phi}'+2\tilde{\phi})\phi_1+\frac{\phi_1^2}{z})\\
        &\quad(1-\frac{\phi_1^2}{8}z^2)+\frac{\phi_1^2}{4z}\\
        &=-\frac{1}{4}(z^3(z\tilde{A}'+3\tilde{A})^2+z(z\tilde{\phi}'+2\tilde{\phi})^2+2(z\tilde{\phi}'+2\tilde{\phi})\phi_1)(1-\frac{\phi_1^2}{8}z^2)\\
        &\quad+\frac{\phi_1^4}{32}z\\
    \end{aligned}
\end{equation}

\begin{equation}
    \begin{aligned}
       V&=-6-phi^2+z^4\tilde{V}_\mathrm{cosh}\\
       V&=-2\phi+z^3\tilde{V}_\mathrm{sinh}
    \end{aligned}
\end{equation}

\begin{equation}
    \begin{aligned}
        \tilde{f}'&=(3+\frac{L^{2}}{2}XV)\frac{1}{z^4}+\frac{\phi_1^2}{8z^2}+\tilde{\xi}_{x}+\mathrm{COM}_s\\
    \end{aligned}
\end{equation}

apparent horizon

\begin{equation}
    \tilde{f}|_h=\textcolor{red}{-}\frac{1}{h^2}\partial_{x'}(h^3\tilde{\xi}+h_{x}e^{-A-\chi})+\frac{h_{x}^2}{h^3}e^{-A-\chi}-\frac{1}{h^3}-\frac{\phi_1^2}{8h}
\end{equation}

\begin{equation}
    \begin{aligned}
        A_0&=\xi'_{x'}-\left(h_xe^{-A-\chi}(A'+\chi')\right)_{x'}\textcolor{red}{+}\frac{f'}{h}\textcolor{red}{+}\frac{h_x^2}{h}e^{-A-\chi}(A'+\chi')\textcolor{red}{-}\frac{f-h_x^2 e^{-A-\chi}}{h^2}\\
        A_1&=\xi'+(e^{-A-\chi})_{x'}+h_x(e^{-A-\chi})'+(\xi+h_x e^{-A-\chi})'\textcolor{red}{-}\frac{2h_x}{h}e^{-A-\chi}\\
        &=(e^{-A-\chi})_{x'}+2(\xi+h_x e^{-A-\chi})'\textcolor{red}{-}\frac{2h_x}{h}e^{-A-\chi}\\
        A_2&=e^{-A-\chi}\\
        B&=(h^3\dot{\tilde{\xi}}-h_xe^{-A-\chi}h^3(\dot{\tilde{A}}+\dot{\tilde{\chi}}))_{x'}\textcolor{red}{+}\frac{h^3\dot{\tilde{f}}+h_x^2 e^{-A-\chi}h^3(\dot{\tilde{A}}+\dot{\tilde{\chi}})}{h}
    \end{aligned}
\end{equation}

\subsubsection{adapted coordinate to fix the horizon}

假设视界位于$z=h(v,x)$处，做坐标变换
\begin{equation}
    \left\{
        \begin{aligned}
            z'&=\frac{z}{h}\\
            v'&=v\\
            x'&=x
        \end{aligned}
    \right.
\end{equation}

则$z=h\Leftrightarrow z'=1$，并且
\begin{equation}
    \left\{
        \begin{aligned}
            dz'&=\frac{dz}{h}-\frac{z}{h^2}(\dot{h}dv+h_xdx)\\
            dv'&=dv\\
            dx'&=dx
        \end{aligned}
    \right.
\end{equation}
\begin{equation}
    \left\{
        \begin{aligned}
            \partial_z&=\frac{1}{h}\partial_{z'}\\
            \partial_v&=\partial_{v'}-\frac{z}{h^2}\dot{h}\partial_{z'}=\partial_{v'}-\frac{z'}{h}\dot{h}\partial_{z'}=\partial_{v'}-z'\dot{h}\partial_z\\
            \partial_x&=\partial_{x'}-\frac{z}{h^2}h_x\partial_{z'}=\partial_{x'}-\frac{z'}{h}h_x\partial_{z'}=\partial_{x'}-z'h_x\partial_z
        \end{aligned}
    \right.
\end{equation}
\begin{equation}
    \left\{
        \begin{aligned}
            \partial_{z'}&=h\partial_z\\
            \partial_{v'}&=\partial_v+\frac{z}{h^2}\dot{h}\partial_{z'}=\partial_v+\frac{z'}{h}\dot{h}\partial_{z'}=\partial_v+z'\dot{h}\partial_z\\
            \partial_{x'}&=\partial_x+\frac{z}{h^2}h_x\partial_{z'}=\partial_x+\frac{z'}{h}h_x\partial_{z'}=\partial_x+z'h_x\partial_z
        \end{aligned}
    \right.
\end{equation}

在解析上用$(v,z,x)$坐标，在数值上用$(v',z',x')$坐标，$z'$方向在[0,1]区间采点，$z'$方向和$x'$方向的求导矩阵分别对应$\partial_{z'}$和$\partial_{x'}$，RK4演化用的时间导数对应$\partial_{v'}$。而所有公式里出现的导数都要用$\partial_v$, $\partial_z$和$\partial_x$。

因为$dz=hdz'+z'\dot{h}dv'+z'h_xdx'$，带撇坐标下的方程，其实等价于用如下度规推出的方程：
\begin{equation}
    ds^2=\frac{L^2}{(z'h)^2}(-[f(v,z,x)e^{-\chi(v,z,x)}-e^{A(v,z,x)}\xi(v,z,x)^2+2z'e^{-\chi}\dot{h}]dv'^2-2e^{-\chi}hdv'dz'-2(\xi e^A+z'e^{-\chi}h_x)dv'dx'+e^Adx'^2+e^{-A}dy^2)
\end{equation}

带撇坐标下的方程为：
\begin{equation}
    \chi_{z'}=\frac{z'}{4}(A_{z'}^2+\phi_{z'}^2)
\end{equation}
\begin{equation}
    \begin{aligned}
        \left(\frac{e^{A+\chi}}{z'^2}\xi_{z'}\right)_{z'}&=-\frac{h}{z'^2}(A+\chi)_{x'z'}-\frac{2h}{z'^3}\chi_{x'}+\frac{h}{z'^2}(A_{x'}A_{z'}+\phi_{x'}\phi_{z'})\\
        &\quad+h_x((A+\chi)_{z'z'}/z'+(A+\chi)_{z'}/z'^2+2\chi_{z'}/z'^2-A_{z'}^2/z'-\phi_{z'}^2/z')\\
        &=\frac{h}{z'^2}[-(A+\chi)_{x'z'}-\frac{2}{z'}\chi_{x'}+(A_{x'}A_{z'}+\phi_{x'}\phi_{z'})]\\
        &\quad+\frac{h_x}{z'}[(A+\chi)_{z'z'}+(A-\chi)_{z'}/z']\\
        \xi_{z'}&=\left(\frac{e^{A+\chi}}{z'^2}\xi_{z'}\right)\frac{z'^2}{e^{A+\chi}}
    \end{aligned}
\end{equation}
\begin{equation}
    \begin{aligned}
        h^2\mathrm{COM}=\frac{e^{A+\chi}}{4}\xi_{z'}^{2}+\frac{h\xi_{x}}{z'}-\frac{h\xi_{xz'}}{2}-\frac{h^2e^{-A-\chi}}{4}(2A_x\chi_x+\chi_{x}^2-\phi_x^2)+\frac{h^2(e^{-A-\chi})_{xx}}{2}+h_x()
    \end{aligned}
\end{equation}
\begin{equation}
    \begin{aligned}
        (\frac{f}{z'^{3}})_{z'}&=\frac{L^{2}}{2z'^{4}}e^{-\chi}V+\frac{h\xi_{x}}{z'^{3}}+\frac{h^2\mathrm{COM}}{z'^{2}}
    \end{aligned}
\end{equation}
\begin{equation}
    \begin{aligned}
        z(\frac{\dot{A}}{z})^{\prime}&=\dot{A}^{\prime}-\frac{\dot{A}}{z}\\
        &=\frac{z^{2}}{2}(\frac{fA^{\prime}-\xi A_{x}}{z^{2}})^{\prime}+\frac{(e^{-A-\chi}A_{x}-\xi A^{\prime})_{x}}{2}+\mathrm{COM}\\
        &=\frac{(fA^{\prime}-\xi A_{x})'}{2}-\frac{fA^{\prime}-\xi A_{x}}{z}+\frac{(e^{-A-\chi}A_{x}-\xi A^{\prime})_{x}}{2}+\mathrm{COM}
    \end{aligned}
\end{equation}
\begin{equation}
    \begin{aligned}
        z(\frac{\dot{\phi}}{z})^{\prime}&=\dot{\phi}^{\prime}-\frac{\dot{\phi}}{z}\\
        &=\frac{z^{2}}{2}(\frac{f\phi'-\xi\phi_{x}}{z^{2}})^{\prime}+\frac{(e^{-A-\chi}\phi_{x}-\xi\phi')_{x}}{2}-\frac{L^{2}}{2z^{2}}e^{-\chi}V_\phi\\
        &=\frac{(f\phi'-\xi\phi_{x})'}{2}-\frac{f\phi'-\xi\phi_{x}}{z}+\frac{(e^{-A-\chi}\phi_{x}-\xi\phi')_{x}}{2}-\frac{L^{2}}{2z^{2}}e^{-\chi}V_\phi
    \end{aligned}
\end{equation}

\subsubsection{adapted coordinate to fix the horizon-2}

假设视界位于$z=H(v,x)$处，做坐标变换
\begin{equation}
    \left\{
        \begin{aligned}
            z'&=\frac{z_0}{H}z\\
            v'&=v\\
            x'&=x
        \end{aligned}
    \right.
\end{equation}

则$z=H\Leftrightarrow z'=z_0$。令$h=\frac{H}{z_0}$，则
\begin{equation}
    \left\{
        \begin{aligned}
            z'&=\frac{z}{h}\\
            v'&=v\\
            x'&=x
        \end{aligned}
    \right.
\end{equation}

并且
\begin{equation}
    \left\{
        \begin{aligned}
            dz'&=\frac{dz}{h}-\frac{z}{h^2}(\dot{h}dv+h_xdx)\\
            dv'&=dv\\
            dx'&=dx
        \end{aligned}
    \right.
\end{equation}
\begin{equation}
    \left\{
        \begin{aligned}
            \partial_z&=\frac{1}{h}\partial_{z'}\\
            \partial_v&=\partial_{v'}-\frac{z}{h^2}\dot{h}\partial_{z'}=\partial_{v'}-\frac{z'}{h}\dot{h}\partial_{z'}=\partial_{v'}-z'\dot{h}\partial_z\\
            \partial_x&=\partial_{x'}-\frac{z}{h^2}h_x\partial_{z'}=\partial_{x'}-\frac{z'}{h}h_x\partial_{z'}=\partial_{x'}-z'h_x\partial_z
        \end{aligned}
    \right.
\end{equation}
\begin{equation}
    \left\{
        \begin{aligned}
            \partial_{z'}&=h\partial_z\\
            \partial_{v'}&=\partial_v+\frac{z}{h^2}\dot{h}\partial_{z'}=\partial_v+\frac{z'}{h}\dot{h}\partial_{z'}=\partial_v+z'\dot{h}\partial_z\\
            \partial_{x'}&=\partial_x+\frac{z}{h^2}h_x\partial_{z'}=\partial_x+\frac{z'}{h}h_x\partial_{z'}=\partial_x+z'h_x\partial_z
        \end{aligned}
    \right.
\end{equation}

在解析上用$(v,z,x)$坐标，在数值上用$(v',z',x')$坐标，$z'$方向在[0,1]区间采点，$z'$方向和$x'$方向的求导矩阵分别对应$\partial_{z'}$和$\partial_{x'}$，RK4演化用的时间导数对应$\partial_{v'}$。而所有公式里出现的导数都要用$\partial_v$, $\partial_z$和$\partial_x$。

与第一种adapted coordinate相比，需要注意的只有在使用视界条件以及求视界的时间演化的时候，此时真正的视界为$H$，并且
\begin{equation}
    \begin{aligned}
        H&=z_0 h\\
        H_x&=z_0 h_x\\
        \dot{H}&=z_0 \dot{h}
    \end{aligned}
\end{equation}

\subsubsection{general adapted coordinate}

假设视界位于$z=h(v,x)$处，做坐标变换
\begin{equation}
    \left\{
        \begin{aligned}
            z'&=F(v,z,x)\\
            v'&=v\\
            x'&=x
        \end{aligned}
    \right.
\end{equation}

则$z=h\Leftrightarrow z'=F(v,h(v,x),x)$，并且
\begin{equation}
    \left\{
        \begin{aligned}
            dz'&=F'dz+\dot{F}dv+F_xdx\\
            dv'&=dv\\
            dx'&=dx
        \end{aligned}
    \right.
\end{equation}
\begin{equation}
    \left\{
        \begin{aligned}
            \partial_z&=F'\partial_{z'}\\
            \partial_v&=\partial_{v'}+\dot{F}\partial_{z'}=\partial_{v'}+\frac{\dot{F}}{F'}\partial_z\\
            \partial_x&=\partial_{x'}+F_x\partial_{z'}=\partial_{x'}+\frac{F_x}{F'}\partial_z
        \end{aligned}
    \right.
\end{equation}
\begin{equation}
    \left\{
        \begin{aligned}
            \partial_{z'}&=\frac{1}{F'}\partial_z\\
            \partial_{v'}&=\partial_v-\dot{F}\partial_{z'}=\partial_v-\frac{\dot{F}}{F'}\partial_z\\
            \partial_{x'}&=\partial_x-F_x\partial_{z'}=\partial_x-\frac{F_x}{F'}\partial_z
        \end{aligned}
    \right.
\end{equation}

在解析上用$(v,z,x)$坐标，在数值上用$(v',z',x')$坐标，$z'$方向在[0,1]区间采点，$z'$方向和$x'$方向的求导矩阵分别对应$\partial_{z'}$和$\partial_{x'}$，RK4演化用的时间导数对应$\partial_{v'}$。而所有公式里出现的导数都要用$\partial_v$, $\partial_z$和$\partial_x$。

\subsubsection{horizon and horizon evolution}

The apparent horizon $z=h(v,x)$:
\begin{equation}\label{h_eq}
    (\partial_{x}+h_{x}\partial_{z})(\xi+h_{x}e^{-A-\chi})=\frac{1}{h}(f-h_{x}^{2}e^{-A-\chi})
\end{equation}

with the fields $(f,\xi,A,\chi)$ understood as functions of $v, z=h(v,x)$ and $x$. Note that the above equation does not involve $\partial_{v}$, which means that the apparent horizon is completely determined by the spatial configuration of $(f,\xi,A,\chi)$ at a given instant $v$. Evolution of the apparent horizon:
\begin{equation}\label{pv_h_eq}
    \begin{aligned}
        &\quad(\partial_{x}+h_{x}\partial_{z})([\partial_{v}+\dot{h}\partial_{z}]\xi+\dot{h}_{x}e^{-A-\chi}+h_{x}[\partial_{v}+\dot{h}\partial_{z}]e^{-A-\chi})+\dot{h}_{x}(\xi+h_{x}e^{-A-\chi})^{\prime}\\
        &=\frac{1}{h}([\partial_{v}+\dot{h}\partial_{z}]f-2h_{x}\dot{h}_{x}e^{-A-\chi}-h_{x}^{2}[\partial_{v}+\dot{h}\partial_{z}]e^{-A-\chi})-h^{-2}\dot{h}(f-h_{x}^{2}e^{-A-\chi})
    \end{aligned}
\end{equation}

在均匀情况退化为
\begin{equation}
    f(v,z=h,x)=f_h=0
\end{equation}
\begin{equation}
    \dot{h}=-\frac{\dot{f}_h}{f'_h}
\end{equation}
\begin{equation}
    \dot{f}_h=-\frac{h}{2}(\dot{A}_h^2+\dot{\phi}_h^2)
\end{equation}
（上式只对视界处成立，因为用到了$f_h=0$来化简）

\eqref{h_eq}即为（注意在视界处$\partial_{x'}=\partial_x+h_x\partial_z$）
\begin{equation}
    \partial_{x'}(\xi+h_{x}e^{-A-\chi})-\frac{1}{h}(f-h_{x}^{2}e^{-A-\chi})=0
\end{equation}

亦即
\begin{equation}
    f=h\partial_{x'}(\xi+h_{x}e^{-A-\chi})+h_{x}^{2}e^{-A-\chi}
\end{equation}


或者展开得到（好像没必要）
\begin{equation}
    \begin{aligned}
        0&=-\frac{f}{h}+\xi_x+h_x\xi'-e^{-A-\chi}[h_x(A_x+h_x A'+\chi_x+h_x\chi'-\frac{h_x}{h})-h_{xx}]\\
        &=-\frac{f}{h}+\xi_{x'}-e^{-A-\chi}[h_x(A_{x'}+\chi_{x'}-\frac{h_x}{h})-h_{xx}]\\
        &=-\frac{f}{h}+\xi_{x'}-e^{-A-\chi}(h_x[(A+\chi)_{x'}-\frac{h_x}{h}]-h_{xx})
    \end{aligned}
\end{equation}

亦即
\begin{equation}
    f=h[\xi_{x'}-e^{-A-\chi}(h_x[(A+\chi)_{x'}-\frac{h_x}{h}]-h_{xx})]
\end{equation}

或者
\begin{equation}
    \partial_{x}(\xi+h_{x}e^{-A-\chi})+h_{x}\partial_{z}(\xi+h_{x}e^{-A-\chi})=\frac{1}{h}(f-h_{x}^{2}e^{-A-\chi})
\end{equation}

用\eqref{h_eq}消去\eqref{pv_h_eq}中的$h_{xx}$，即\eqref{pv_h_eq}+$(\partial_v+\dot{h}\partial_z)(A+\chi)$\eqref{h_eq}，得到
\begin{equation}
    A_0\dot{h}+A_1\dot{h}_x+A_2\dot{h}_{xx}+B=0
\end{equation}

其中
\begin{equation}
    \begin{aligned}
        A_0&=(\frac{f}{h}-f')\frac{1}{h}+(A'+\chi')(\xi_x+h_x\xi'-\frac{f}{h})+\xi'_x+h_x\xi''-e^{-A-\chi}h_x[(A'+\chi')_x+h_x(A'+\chi')'+\frac{h_x}{h^{2}}]\\
        &=(\frac{f}{h}-f')\frac{1}{h}+(A'+\chi')(\xi_{x'}-\frac{f}{h})+\xi'_{x'}-e^{-A-\chi}h_x[(A'+\chi')_{x'}+\frac{h_x}{h^{2}}]
    \end{aligned}
\end{equation}
\begin{equation}
    \begin{aligned}
        A_1&=2\xi'-e^{-A-\chi}(A_x+3h_x A'+\chi_x+3h_x \chi'-\frac{2h_x}{h})\\
        &=2\xi'-e^{-A-\chi}(A_{x'}+\chi_{x'}+2h_x A'+2h_x \chi'-\frac{2h_x}{h})\\
        &=2\xi'-e^{-A-\chi}[(A+\chi)_{x'}+2h_x(A'+\chi'-\frac{1}{h})]
    \end{aligned}
\end{equation}
\begin{equation}
    A_2=e^{-A-\chi}
\end{equation}
\begin{equation}
    \begin{aligned}
        B&=-\frac{\dot{f}}{h}+(\dot{A}+\dot{\chi})(\xi_x+h_x\xi'-\frac{f}{h})+\dot{\xi}_x+h_x\dot{\xi}'-e^{-A-\chi}h_x[(\dot{A}+\dot{\chi})_x+h_x(\dot{A}+\dot{\chi})']\\
        &=-\frac{\dot{f}}{h}+(\dot{A}+\dot{\chi})(\xi_{x'}-\frac{f}{h})+\dot{\xi}_{x'}-e^{-A-\chi}h_x(\dot{A}+\dot{\chi})_{x'}
    \end{aligned}
\end{equation}

也可以直接从\eqref{pv_h_eq}+$(\partial_v+\dot{h}\partial_z)(A+\chi)$\eqref{h_eq}读出
\begin{equation}
    \begin{aligned}
        A_0&=\xi'_{x'}+\left(h_x(e^{-A-\chi})'\right)_{x'}-\frac{f'}{h}+\frac{h_x^2}{h}(e^{-A-\chi})'+\frac{f-h_x^2 e^{-A-\chi}}{h^2}\\
        &\quad+(A+\chi)'[\partial_{x'}(\xi+h_{x}e^{-A-\chi})-\frac{1}{h}(f-h_{x}^{2}e^{-A-\chi})]\\
        &=\xi'_{x'}-(A'+\chi')_{x'}h_x e^{-A-\chi}-\frac{f'}{h}+\frac{f-h_x^2 e^{-A-\chi}}{h^2}\\
        &\quad+(A'+\chi')(\xi_{x'}-\frac{f}{h})\\
    \end{aligned}
\end{equation}
\begin{equation}
    \begin{aligned}
        A_1&=\xi'+(e^{-A-\chi})_{x'}+h_x(e^{-A-\chi})'+(\xi+h_x e^{-A-\chi})'+\frac{2h_x}{h}e^{-A-\chi}\\
        &=(e^{-A-\chi})_{x'}+2(\xi+h_x e^{-A-\chi})'+\frac{2h_x}{h}e^{-A-\chi}
    \end{aligned}
\end{equation}
\begin{equation}
    A_2=e^{-A-\chi}
\end{equation}
\begin{equation}
    \begin{aligned}
        B&=(\dot{\xi}+h_x\partial_v(e^{-A-\chi}))_{x'}-\frac{\dot{f}-h_x^2\partial_v(e^{-A-\chi})}{h}\\
        &\quad+(\dot{A}+\dot{\chi})[\partial_{x'}(\xi+h_{x}e^{-A-\chi})-\frac{1}{h}(f-h_{x}^{2}e^{-A-\chi})]\\
        &=\dot{\xi}_{x'}-(\dot{A}+\dot{\chi})_{x'}h_xe^{-A-\chi}-\frac{\dot{f}}{h}\\
        &\quad+(\dot{A}+\dot{\chi})(\xi_{x'}-\frac{f}{h})
    \end{aligned}
\end{equation}

如果直接用\eqref{pv_h_eq}，则
\begin{equation}
    \begin{aligned}
        A_0&=\xi'_{x'}+\left(h_x(e^{-A-\chi})'\right)_{x'}-\frac{f'}{h}+\frac{h_x^2}{h}(e^{-A-\chi})'+\frac{f-h_x^2 e^{-A-\chi}}{h^2}\\
        A_1&=\xi'+(e^{-A-\chi})_{x'}+h_x(e^{-A-\chi})'+(\xi+h_x e^{-A-\chi})'+\frac{2h_x}{h}e^{-A-\chi}\\
        &=(e^{-A-\chi})_{x'}+2(\xi+h_x e^{-A-\chi})'+\frac{2h_x}{h}e^{-A-\chi}\\
        A_2&=e^{-A-\chi}\\
        B&=(\dot{\xi}+h_x\partial_v(e^{-A-\chi}))_{x'}-\frac{\dot{f}-h_x^2\partial_v(e^{-A-\chi})}{h}
    \end{aligned}
\end{equation}

\subsubsection{apparent horizon}

第一种理解：

in time-slice $v=v_0$, apparent horizon $\Phi(z,x):=z-h(x)$, the normal vector
\begin{equation}
    k_i\sim\partial_i\Phi=(1,-h_x,0)
\end{equation}

上式中的法矢量是在time-slice $v=v_0$上定义的，把它延拓到四维时空
\begin{equation}
    k_\mu=\lambda(v,z,x,y)(k_0,1,-h_x,0)
\end{equation}

if $k_\mu$ is null
\begin{equation}
    k^\mu k_\mu=0\Rightarrow k_0=\xi h_x + \frac{f}{2} + \frac{e^{- A - \chi} h_x^2}{2}
\end{equation}

第二种理解：

\textcolor{red}{这种理解是错的！这相当于要求了所有时刻的表观视界堆叠起来的余1维超曲面是类光的，但它应该是类空的。}

assuming the hypersurface of apparent horizon is $\Phi(z,x):=z-h(v,x)$, so the normal vector (or the congruence may be defined as)
\begin{equation}
    k_\mu\sim\partial_\mu\Phi=(-\dot{h},1,-h_x,0)
\end{equation}

let
\begin{equation}
    k_\mu=\lambda(v,z,x,y)\partial_\mu\Phi=\lambda(v,z,x,y)(-\dot{h},1,-h_x,0)
\end{equation}

Requiring that $k$ be null fixes the time derivative $\partial_v\Phi$ in terms of spatial derivatives of $\Phi$, such that
\begin{equation}
    k^\mu k_\mu=0\Rightarrow\dot{h}=- \xi h_x - \frac{f}{2} - \frac{e^{- A - \chi} h_x^2}{2}
\end{equation}

用这种理解来进行后续推导的话，得到的表观视界条件相同，而测地线条件中会出现$\dot{h}_x$项，如果用上式得到的$\dot{h}$来化简会得到
\begin{equation}
    0=2 \xi' h_x + f' - e^{-A-\chi} A' h_x^2 - e^{-A-\chi} \chi' h_x^2
\end{equation}

方法1：

Requiring the congruence to satisfy the (affinely parameterized) geodesic equation $k^\alpha k_{\beta;\alpha}=0$ determines the time derivative of the multiplier function $\lambda$ in terms of its spatial derivatives.

$v$分量由$z$和$x$分量组合而来，$z$分量给出$\dot{\lambda}$，而$x$分量给出测地线条件。
\begin{equation}
    \dot{\lambda}=\lambda \xi' h_x + \frac{\lambda f'}{2} - \frac{\lambda e^{-A-\chi} A' h_x^2}{2} - \frac{\lambda e^{-A-\chi} \chi' h_x^2}{2} - \xi \lambda_x + \frac{f \lambda'}{2} - e^{-A-\chi} \lambda_x h_x - \frac{e^{-A-\chi} \lambda' h_x^2}{2}
\end{equation}

and the geodesic condition
\begin{equation}
    0=\xi h_{xx} + \xi_x h_x + \xi' h_x^2 + \frac{f_x}{2} + \frac{f' h_x}{2} + \dot{h}_x - \frac{e^{-A-\chi} A_x h_x^2}{2} - \frac{e^{-A-\chi} A' h_x^3}{2} - \frac{e^{-A-\chi} \chi_x h_x^2}{2} - \frac{e^{-A-\chi} \chi' h_x^3}{2} + e^{-A-\chi} h_x h_{xx}
\end{equation}

Given these time derivatives, one may then compute the expansion via $\theta=\nabla\cdot k$.
\begin{equation}
    \begin{aligned}
        0&=\theta=\nabla_\mu k^\mu\Rightarrow\\
        0&=\xi_x + \xi' h_x - e^{-A-\chi} A_x h_x - e^{-A-\chi} A' h_x^2 - e^{-A-\chi} \chi_x h_x - e^{-A-\chi} \chi' h_x^2 + e^{-A-\chi} h_{xx} + \frac{f}{h} - \frac{e^{-A-\chi} h_x^2}{h}
    \end{aligned}
\end{equation}

i.e.
\begin{equation}
    (\partial_{x}+h_{x}\partial_{z})(\xi+h_{x}e^{-A-\chi})=\textcolor{red}{-}\frac{1}{h}(f-h_{x}^{2}e^{-A-\chi})
\end{equation}

方法2

取$\lambda=1$，则测地线方程应该有最一般的形式，而不一定是仿射参数仿射参数下的测地线方程（如果适当的选取$\lambda$，可以使得仿射参数下的测地线方程$k^\mu\nabla_\mu k_\nu=0$成立），则
\begin{equation}
    k^\mu\nabla_\mu k_\nu-\kappa k_\nu=0
\end{equation}

expansion
\begin{equation}
    \hat{\theta}=\theta-\kappa=\nabla_\mu k^\mu-\kappa=0
\end{equation}

so
\begin{equation}
    k^\mu\nabla_\mu k_\nu-(\nabla_\mu k^\mu)k_\nu=0
\end{equation}

$v$分量由$z$和$x$分量组合而来，$z$分量给出表观视界条件，而$x$分量给出测地线条件。得到的结果和方法1一样（注意测地线条件可以利用表观视界条件来化简）。

总结：

测地线条件
\begin{equation}
    \begin{aligned}
        0&=2 \xi h_{xx} + \xi_x h_x + \xi' h_x^2 + f_x + f' h_x + e^{-A-\chi} h_x h_{xx} - \frac{f h_x}{h} + \frac{e^{-A-\chi} h_x^3}{h}\\
        &=(\partial_x+h_x\partial_z)(\xi h_x+f) + \xi h_{xx} + e^{-A-\chi} h_x h_{xx} - \frac{f h_x}{h} + \frac{e^{-A-\chi} h_x^3}{h}\\
    \end{aligned}
\end{equation}

表观视界条件
\begin{equation}
    (\partial_{x}+h_{x}\partial_{z})(\xi+h_{x}e^{-A-\chi})=\textcolor{red}{-}\frac{1}{h}(f-h_{x}^{2}e^{-A-\chi})
\end{equation}

\begin{equation}
    f|_{h}=\textcolor{red}{-}h\partial_{x'}(\xi+h_{x}e^{-A-\chi})+h_{x}^{2}e^{-A-\chi}
\end{equation}

\begin{equation}
    \xi_x + \xi' h_x - e^{-A-\chi} A_x h_x - e^{-A-\chi} A' h_x^2 - e^{-A-\chi} \chi_x h_x - e^{-A-\chi} \chi' h_x^2 + e^{-A-\chi} h_{xx} + \frac{f}{h} - \frac{e^{-A-\chi} h_x^2}{h}=0
\end{equation}

对$h(x)$变分
\begin{equation}
    \begin{aligned}
        0&=(\partial_{x}+h_{x}\partial_{z})\delta(\xi+h_{x}e^{-A-\chi})+(\delta h_{x})\partial_{z}(\xi+h_{x}e^{-A-\chi})+\delta[\frac{1}{h}(f-h_{x}^{2}e^{-A-\chi})]\\
        &=(\partial_{x}+h_{x}\partial_{z})(\xi'\delta h+(\delta h_{x})e^{-A-\chi}+h_{x}(e^{-A-\chi})'\delta h)+(\delta h_{x})\partial_{z}(\xi+h_{x}e^{-A-\chi})\\
        &\quad+\delta(\frac{1}{h})(f-h_{x}^{2}e^{-A-\chi})+\frac{1}{h}\delta(f-h_{x}^{2}e^{-A-\chi})\\
        &=(\partial_{x}+h_{x}\partial_{z})(\xi'\delta h+(\delta h_{x})e^{-A-\chi}+h_{x}(e^{-A-\chi})'\delta h)+(\delta h_{x})\partial_{z}(\xi+h_{x}e^{-A-\chi})\\
        &\quad-\frac{1}{h^2}(f-h_{x}^{2}e^{-A-\chi})\delta h+\frac{1}{h}(f'\delta h-2h_{x}e^{-A-\chi}\delta h_x-h_{x}^{2}(e^{-A-\chi})'\delta h)\\
        &=\partial_{x}[\xi'\delta h+(\delta h_{x})e^{-A-\chi}+h_{x}(e^{-A-\chi})'\delta h]+h_{x}\partial_{z}[\xi'\delta h+(\delta h_{x})e^{-A-\chi}+h_{x}(e^{-A-\chi})'\delta h]+(\delta h_{x})\partial_{z}(\xi+h_{x}e^{-A-\chi})\\
        &\quad-\frac{1}{h^2}(f-h_{x}^{2}e^{-A-\chi})\delta h+\frac{1}{h}(f'\delta h-2h_{x}e^{-A-\chi}\delta h_x-h_{x}^{2}(e^{-A-\chi})'\delta h)\\
        &=\delta h[(\partial_{x}+h_{x}\partial_{z})(\xi'+h_{x}(e^{-A-\chi})')-\frac{1}{h^2}(f-h_{x}^{2}e^{-A-\chi})+\frac{1}{h}(f'-h_{x}^{2}(e^{-A-\chi})')]\\
        &\quad+\delta h_x[(\xi'+(e^{-A-\chi})_x+h_{x}(e^{-A-\chi})')+h_{x}(e^{-A-\chi})'+\partial_{z}(\xi+h_{x}e^{-A-\chi})-\frac{2h_{x}e^{-A-\chi}}{h}]\\
        &\quad+\delta h_{xx}e^{-A-\chi}
    \end{aligned}
\end{equation}

\begin{equation}
    V_0\delta h+V_1\delta h_x+V_2\delta h_{xx}=0
\end{equation}

\begin{equation}
    \begin{aligned}
        V_0&=- \frac{f}{h^{2}} + \frac{e^{-A-\chi} h_x^{2}}{h^{2}}\\
        V_1&=\xi' - e^{-A-\chi} A_x - 2 e^{-A-\chi} A' h_x - e^{-A-\chi} \chi_x - 2 e^{-A-\chi} \chi' h_x - \frac{2 e^{-A-\chi} h_x}{h}\\
        &=\xi' + (e^{-A-\chi})_x + 2 (e^{-A-\chi})' h_x - \frac{2 e^{-A-\chi} h_x}{h}\\
        V_2&=e^{-A-\chi}\\
    \end{aligned}
\end{equation}

视界的演化
\begin{equation}
    \begin{aligned}
        &\quad(\partial_{x}+h_{x}\partial_{z})([\partial_{v}+\dot{h}\partial_{z}]\xi+\dot{h}_{x}e^{-A-\chi}+h_{x}[\partial_{v}+\dot{h}\partial_{z}]e^{-A-\chi})+\dot{h}_{x}(\xi+h_{x}e^{-A-\chi})^{\prime}\\
        &=\textcolor{red}{-}\frac{1}{h}([\partial_{v}+\dot{h}\partial_{z}]f-2h_{x}\dot{h}_{x}e^{-A-\chi}-h_{x}^{2}[\partial_{v}+\dot{h}\partial_{z}]e^{-A-\chi})\textcolor{red}{+}h^{-2}\dot{h}(f-h_{x}^{2}e^{-A-\chi})
    \end{aligned}
\end{equation}

\begin{equation}
    A_0\dot{h}+A_1\dot{h}_x+A_2\dot{h}_{xx}+B=0
\end{equation}

\begin{equation}
    \begin{aligned}
        A_0&=\xi'_{x'}+\left(h_x(e^{-A-\chi})'\right)_{x'}\textcolor{red}{+}\frac{f'}{h}\textcolor{red}{-}\frac{h_x^2}{h}(e^{-A-\chi})'\textcolor{red}{-}\frac{f-h_x^2 e^{-A-\chi}}{h^2}\\
        &=(\xi+h_xe^{-A-\chi})'_{x'}\textcolor{red}{+}\frac{f'}{h}\textcolor{red}{-}\frac{h_x^2}{h}(e^{-A-\chi})'\textcolor{red}{-}\frac{f-h_x^2 e^{-A-\chi}}{h^2}\\
        A_1&=\xi'+(e^{-A-\chi})_{x'}+h_x(e^{-A-\chi})'+(\xi+h_x e^{-A-\chi})'\textcolor{red}{-}\frac{2h_x}{h}e^{-A-\chi}\\
        &=(e^{-A-\chi})_{x'}+2(\xi+h_x e^{-A-\chi})'\textcolor{red}{-}\frac{2h_x}{h}e^{-A-\chi}\\
        A_2&=e^{-A-\chi}\\
        B&=(\dot{\xi}+h_x\partial_v(e^{-A-\chi}))_{x'}\textcolor{red}{+}\frac{\dot{f}-h_x^2\partial_v(e^{-A-\chi})}{h}
    \end{aligned}
\end{equation}

\subsubsection{evolution procedure}

初值：$A$,$\phi$,$\xi_3$,$f_3$

1. $A,\phi\rightarrow\chi'\rightarrow\chi$，边界条件为$\chi|_B=0$

2. $A,\phi,\chi\rightarrow\left(\frac{e^{A+\chi}}{z^2}\xi'\right)'\rightarrow\left(\frac{e^{A+\chi}}{z^2}\xi'\right)$，边界条件为$\left(\frac{e^{A+\chi}}{z^2}\xi'\right)|_B=\frac{e^{A+\chi}\xi'''}{2}|_B=3e^{A+\chi}|_B\xi_3$

$\left(\frac{e^{A+\chi}}{z^2}\xi'\right)\rightarrow\xi'\rightarrow\xi$，边界条件为$\xi|_B=0$

3. $A,\phi,\chi,\xi\rightarrow\tilde{f}'\rightarrow\tilde{f}$，边界条件为$\tilde{f}'|_B=f_3$

$\tilde{f}\rightarrow f=1+z^2\tilde{f}$
\begin{equation}
    \begin{aligned}
        (\frac{1+z^2\tilde{f}}{z^{3}})^{\prime}&=\frac{\xi_{x}}{z^{3}}+\frac{L^{2}}{2z^{4}}e^{-\chi}V+\frac{\mathrm{COM}}{z^{2}}\\
        z(\frac{\tilde{f}}{z})^{\prime}&=\tilde{f}'-\frac{\tilde
        {f}}{z}\\
        &=\frac{\xi_{x}}{z^2}+\frac{1}{z^3}(3+\frac{L^{2}}{2}e^{-\chi}V)+\frac{\mathrm{COM}}{z}
    \end{aligned}
\end{equation}

$3'$. $A,\phi,\chi,\xi\rightarrow f'\rightarrow f$，边界条件为$f|_H=(h\partial_{x'}(\xi+h_{x}e^{-A-\chi})+h_{x}^{2}e^{-A-\chi})|_H$
\begin{equation}
    \begin{aligned}
        (\frac{f}{z^{3}})^{\prime}&=\frac{L^{2}}{2z^{4}}e^{-\chi}V+\frac{\xi_{x}}{z^{3}}+\frac{\mathrm{COM}}{z^{2}}\\
        zf'-3f&=\frac{L^{2}}{2}e^{-\chi}V+z\xi_{x}+z^2\mathrm{COM}\\
        &=\frac{L^{2}}{2}e^{-\chi}V+2z\xi_{x}+z^2(\mathrm{COM}-\frac{\xi_x}{z})\\
        &=\frac{L^{2}}{2}e^{-\chi}V+2z\xi_{x}+z^2\mathrm{COM}_f\\
        \mathrm{COM}_f&=\mathrm{COM}-\frac{\xi_x}{z}\\
        &=\frac{e^{A+\chi}}{4}\xi^{\prime2}-\frac{\xi_{x}^{\prime}}{2}-\frac{e^{-A-\chi}}{4}(2A_x\chi_x+\chi_{x}^2-\phi_x^2)+\frac{(e^{-A-\chi})_{xx}}{2}
    \end{aligned}
\end{equation}

4. $A,\phi,\chi,\xi,f\rightarrow \dot{A}'\rightarrow \dot{A}$，边界条件为$\dot{A}'|_B=0$

5. $A,\phi,\chi,\xi,f\rightarrow \dot{\phi}'\rightarrow \dot{\phi}$，边界条件为$\dot{\phi}'|_B=0$

如果要固定视界，还需要提供$h$的初值，并且需要计算$\dot{h}$：

6.$\rightarrow\dot{\chi}\rightarrow\dot{\xi}\rightarrow\dot{f}\rightarrow\dot{h}$

\subsubsection{evolution result}

\textbf{跟演化有关的因素:}

\begin{itemize}
    \item 初值的选取：优先考虑视界在1的初值，因此选了source=1.5，具有spinodal不稳定性
    \item 微扰形式：在$x$方向周期为$l$
    
    类似于高斯波包的微扰，可以保证各个模式的微扰都有：
    \begin{equation}
        \delta\phi=0.1z^2(1-z^2)\exp[-10\cos^2(\frac{x}{l}\pi)]
    \end{equation}
    
    正弦微扰：
    \begin{equation}
        \delta\phi=0.1z^2(1-z^2)\sin(\frac{2x}{l}\pi)
    \end{equation}
    
    微扰形式会影响末态构型，不过对程序稳定性影响不大
    \item 径向方向采点个数$N$
    \item $x$方向区间长度$l$
    \item $x$方向采点个数$M$
    
    总的来说，$x$方向网格越细，混叠现象越严重，锯齿越明显，程序崩得越快

    可能不一定，见后面
    \item 时间步长$dt$：满足CFL条件
    \item 方程上的细节
    \begin{itemize}
        \item 提到高阶：测试了$A$和$\phi$再提一阶，让导数算符里没有发散，不过常数项部分仍然有发散，影响不大
        \item 分开算（symbol转numeric）还是合起来算：之前测试会导致崩与不崩，换了初值之后又没复现出来
        \item 上面两条可能与提到高阶有关，最新测试，把所有场都提到最高阶，不固定演化会崩，视界往外走的地方会长出尖，猜测是因为没有涵盖视界导致的崩溃。而之前的代码能运行是因为运气好，两种误差刚好抵消，勉强能达到末态，但其实误差很大
        \item 对$f$求导用$(1+z^2\tilde{f})'$还是$2z\tilde{f}+z^2\tilde{f}'$：之前测试会导致不崩与崩，换了初值之后又没复现出来
        \item $x$方向的二阶求导矩阵：用解析表达式还是一阶求导矩阵的平方：之前测试会导致崩与不崩，换了初值之后又没复现出来
        
        按理来说，对于奇数采点，两者相等；对于偶数采点，两者相差一个正负交替的常数，但是作用到函数之后得到的结果相同（跟函数的光滑性有关）
    \end{itemize}
    \item 是否固定视界：不管是物理上还是数值上，都是合理的，而且从均匀超辐射的结果来看，固定视界确实可以提高稳定性并降低误差，而且使用视界边界条件效果是最好的，可以跑到q=100
    \item $f$用视界边界条件还是能量边界条件：从均匀超辐射的结果来看，使用视界边界条件效果是最好的，可以跑到q=100
\end{itemize}

\textbf{均匀情况测试：$\phi[0]=1.5z$，给定能量，动态收敛到一个静态解}

\begin{itemize}
    \item 不固定视界+能量边界条件：work
    \item 固定视界+能量边界条件：work，不过如果收敛之后还一直跑的话，混叠误差也会积累起来
    \item 固定视界+视界边界条件：崩溃
    
    把$x$方向关了（也就是用退化的公式）可以work

    崩的时候看$h$和其他几个场，都有明显的锯齿，哪怕是本来应该恒等于零的$A$和$\xi$，也出现了比较大的锯齿

    $x$方向网格越细，崩得越快（如果取$M=1$就可以work，不过这等价于是直接把$x$方向关了）

    和$M$的关系好像没那么直接，比较小的时候要看情况，而比较大之后，似乎和他无关了

    不过空间长度确实是越长就崩得越慢（但其实它发生相分离的时间也延后了）

    $N$越大崩的也越快，但是应该没那么敏感
    
    可以确保视界处的边界条件没错，而且换成f=0也没有改善（会缓解），猜测是pvf有问题，因为用-pvf/pzf是崩的最快的，用矩阵求逆求pvh反而没那么快崩。但是把f[0]=0和pvf/pzf一起用，就可以work

    计算结果和symbol转numeric的结果一样，基本排除是公式输错

    排除了是h operator取逆的问题，它确实是可逆的，一直都满足相乘等于单位矩阵的条件，说明周期性边界条件应该够了，不需要再添加边界条件

    关掉pvh，work

    关掉pvA，可以缓解

    关掉xi，即令xi一直=0，反而加剧

    同时关掉A和xi，work
\end{itemize}

\textbf{相分离测试：}

\begin{itemize}
    \item 不固定视界+能量边界条件：某些参数下可以演化到稳定末态，但是误差是O(0.1)。得到的结果f还算光滑，但是对x求导就有锯齿了。对N很敏感，测试结果显示，总体来说N取奇数就会崩，N取偶数就可以
    
    相分离发生时误差最大，在0.4左右，末态时误差在0.1的量级
    在$z$方向，误差都是越原理边界越大，猜测是视界以内的部分误差大。而且在某一个$x$处相对更大一些，猜测是那个点的视界出现在了计算域的外面
    在$x$方向呈锯齿分布，猜测是混叠现象导致的，不过在这种方案下还没占主导地位，不至于让程序崩掉
    
    f3明明加了边界条件，误差也很大，可能是三阶导数本来就不精确

    径向方向采点个数$N$：随N变化不收敛，有的会崩有的不会崩，总体来说是偶数不崩，奇数崩

    $x$方向区间长度$l$和$x$方向采点个数$M$：网格太细的话，混叠越多，锯齿越明显，甚至崩掉；网格太粗的话，函数也会不怎么连续
    \item 固定视界+能量边界条件：大概是在发生相分离的时候崩
    
    锯齿，混叠
    \item 固定视界+视界边界条件：崩得更快
    
    锯齿，混叠
\end{itemize}

\textbf{其它的一些尝试：}

\begin{itemize}
    \item 滤波：在每一个时间步对最终结果滤波，会缓解，崩得会慢一点。砍掉1/2效果又要比砍掉1/3好
    \item 用第二套方案，全用演化方程：均匀构型也崩得很快，说明混叠误差积累主要发生在求pvchi,pvxi,pvf,pvh的过程中，而用约束方程的话不会积累太大的误差
    \begin{itemize}
        \item chi单独用约束方程还是演化方程，基本没变化
        \item xi单独用约束方程，会好一点
        \item f单独用约束方程，会好一点
        \item xi和f都用约束方程，work
        \item 用约束方程求pvxi，work。说明问题出在pvxi
        \item 用约束方程求pvf，有问题。进一步说明问题出在pvxi
    \end{itemize}
    \item 每五个时间步就能用五点差分得到一个pvprimechi，然后和pvchi相减，就能得到pvh，每两步用一次这个方法修正pvh：没有改善，而且两种方法得到的pvh相差比较多，但从公式上看两者似乎是等价的
    \item 在$[0,\frac{2}{1+\cos(\frac{10}\pi{N-1})}]$采点，第11个点恰好是1，然后把视界固定在1
    
    没有太大影响，崩的时间也基本没变化
    \item 把$z'=1$固定成$f$的零点，而不是表观视界：不管是能量边界条件还是视界边界条件，对均匀构型可以work，对相分离情况会崩。视界边界条件崩得会慢很多，而能量边界条件崩得更快了。
    \item 通过对xi和f的约束方程求时间导数来求pvxi和pvf，有很大改善，第二套演化方案也可以work很好了
    \item 把五个场都提到了最高阶，除了标量势外，所有发散都消了，运行的时候误差似乎小了，但不管N取什么值都出现了之前N取奇数出现的尖，导致崩溃。xi的演化方程求的pvxi仍然会让演化很快崩溃
    \item xi和f的约束方程来求pvxi和pvf之后，把$z'=1$固定成$f$的零点，表现得不错，崩的时候都没出现锯齿，很光滑，这个崩溃可能是没覆盖住视界造成的，而不是混叠。但是如果没覆盖住视界再强行演化几步（能量边界条件会比视界边界条件多演化几步），就会出现严重的锯齿
    \item 在此基础上把采点域扩大，似乎有一点变化，但崩的时间更快一点
    \item 在xi或f的两边都施加边界条件，会崩得很快
\end{itemize}

\begin{itemize}
    \item 只要用到演化方程求出的pvxi，不管是用来求演化xi还是演化h，就会有积累得很快的误差
\end{itemize}

\subsubsection{PPT}

假设视界位于$z=h(v,x)$处，做坐标变换
\begin{equation}
    \left\{
        \begin{aligned}
            z'&=\frac{z}{h}\\
            v'&=v\\
            x'&=x
        \end{aligned}
    \right.
    \Rightarrow
    \left\{
        \begin{aligned}
            \partial_z&=\frac{1}{h}\partial_z'\\
            \partial_v&=\partial_{v'}-z'\dot{h}\partial_z\\
            \partial_x&=\partial_{x'}-z'h_x\partial_z
        \end{aligned}
    \right.
\end{equation}

则$z=h\Leftrightarrow z'=1$，并且
\begin{equation}
    \left\{
        \begin{aligned}
            \partial_z&=\frac{1}{h}\partial_z'\\
            \partial_v&=\partial_{v'}-z'\dot{h}\partial_z\\
            \partial_x&=\partial_{x'}-z'h_x\partial_z
        \end{aligned}
    \right.
\end{equation}

视界位置
\begin{equation}
    (\partial_{x}+h_{x}\partial_{z})(\xi+h_{x}e^{-A-\chi})=\frac{1}{h}(f-h_{x}^{2}e^{-A-\chi})
\end{equation}

\begin{equation}
    \begin{aligned}
        &\quad(\partial_{x}+h_{x}\partial_{z})([\partial_{v}+\dot{h}\partial_{z}]\xi+\dot{h}_{x}e^{-A-\chi}+h_{x}[\partial_{v}+\dot{h}\partial_{z}]e^{-A-\chi})+\dot{h}_{x}(\xi+h_{x}e^{-A-\chi})^{\prime}\\
        &=\frac{1}{h}([\partial_{v}+\dot{h}\partial_{z}]f-2h_{x}\dot{h}_{x}e^{-A-\chi}-h_{x}^{2}[\partial_{v}+\dot{h}\partial_{z}]e^{-A-\chi})-h^{-2}\dot{h}(f-h_{x}^{2}e^{-A-\chi})
    \end{aligned}
\end{equation}

\begin{equation}
    A_0\dot{h}+A_1\dot{h}_x+A_2\dot{h}_{xx}+B=0
\end{equation}

\begin{equation}
    \begin{aligned}
        A_0&=(\frac{f}{h}-f')\frac{1}{h}+(A'+\chi')(\xi_{x'}-\frac{f}{h})+\xi'_{x'}-e^{-A-\chi}h_x[(A'+\chi')_{x'}+\frac{h_x}{h^{2}}]\\
        A_1&=(e^{-A-\chi})_{x'}+2(\xi+h_x e^{-A-\chi})'+\frac{2h_x}{h}e^{-A-\chi}\\
        A_2&=e^{-A-\chi}\\
        B&=-\frac{\dot{f}}{h}+(\dot{A}+\dot{\chi})(\xi_{x'}-\frac{f}{h})+\dot{\xi}_{x'}-e^{-A-\chi}h_x(\dot{A}+\dot{\chi})_{x'}
    \end{aligned}
\end{equation}

\begin{equation}
    f|_h=h\partial_{x'}(\xi+h_{x}e^{-A-\chi})+h_{x}^{2}e^{-A-\chi}
\end{equation}

\subsection{inhomogeneous case with spherical topology}

\subsubsection{equation}

The model is same with the case in planar topology, and the BS coordinate
\begin{equation}
    ds^2=\frac{L^2}{z^2}(-[f(v,z,\theta)e^{-\chi(v,z,\theta)}-e^{A(v,z,\theta)}\xi(v,z,\theta)^2]dv^2-2e^{-\chi}dvdz-2\xi e^Advd\theta+e^Ad\theta^2+e^{-A}\sin^2\theta d\varphi^2)
\end{equation}

$E_{zz}$ gives the constraint equation of $\chi$:
\begin{equation}
    \chi'=\frac{z}{4}(A'^2+\phi'^2)
\end{equation}

$E_{z\theta}$ gives the constraint equation of $\xi$:
\begin{equation}
    \begin{aligned}
        \left(\frac{e^{A+\chi}}{z^2}\xi'\right)'&=-\frac{1}{z^{2}}[(A+\chi)^{\prime}_{\theta}+\frac{2\chi_{\theta}}{z}-(A^{\prime}A_{\theta}+\phi'\phi_{\theta})\textcolor{red}{+2\cot\theta A'}]\\
        \xi'&=\left(\frac{e^{A+\chi}}{z^2}\xi'\right)\frac{z^2}{e^{A+\chi}}
    \end{aligned}
\end{equation}

$E_{vz}-\frac{f}{2}E_{zz}+\xi E_{zx}$ gives the constraint equation of $f$:
\begin{equation}
    \begin{aligned}
        \mathrm{COM}&=\frac{e^{A+\chi}}{4}\xi^{\prime2}+\frac{\xi_{\theta}}{z}-\frac{\xi_{\theta}^{\prime}}{2}-\frac{e^{-A-\chi}}{4}(2A_x\chi_x+\chi_{\theta}^2-\phi_x^2)+\frac{(e^{-A-\chi})_{xx}}{2}\\
        &=\frac{e^{A+\chi}}{4}\xi^{\prime2}+\frac{\xi_{\theta}}{z}-\frac{\xi_{\theta}^{\prime}}{2}+\frac{e^{-A-\chi}}{2}[A_x^2+A_x\chi_x+\frac{1}{2}\chi_{\theta}^2+\frac{1}{2}\phi_x^2-(A_{xx}+\chi_{xx})]
    \end{aligned}
\end{equation}
\begin{equation}
    \textcolor{red}{\mathrm{COM}_\theta=\frac{\xi'}{2}-\frac{\xi}{z}+\frac{e^{-A-\chi}\chi_\theta}{2}}
\end{equation}
\begin{equation}
    \begin{aligned}
        (\frac{f}{z^{3}})^{\prime}&=\frac{\xi_{\theta}}{z^{3}}+\frac{L^{2}}{2z^{4}}e^{-\chi}V+\frac{\mathrm{COM}}{z^{2}}\textcolor{red}{-\cot\theta(\frac{\xi'}{2z^2}+\frac{3e^{-A-\chi}A_\theta}{2z^2}+\frac{e^{-A-\chi}\chi_\theta}{2z^2}-\frac{2\xi}{z^3})-\frac{e^{-A-\chi}}{z^2}}\\
        &=\frac{\xi_{\theta}}{z^{3}}+\frac{L^{2}}{2z^{4}}e^{-\chi}V+\frac{\mathrm{COM}\textcolor{red}{-\cot\theta(\mathrm{COM}_\theta-\frac{\xi}{z}+\frac{3e^{-A-\chi}A_\theta}{2})-e^{-A-\chi}}}{z^{2}}\\
        &=(\frac{1+z^2\tilde{f}}{z^{3}})^{\prime}\\
        z(\frac{\tilde{f}}{z})^{\prime}&=\tilde{f}'-\frac{\tilde
        {f}}{z}\\
        &=\frac{\xi_{\theta}}{z^2}+\frac{1}{z^3}(3+\frac{L^{2}}{2}e^{-\chi}V)+\frac{\mathrm{COM}\textcolor{red}{-\cot\theta(\mathrm{COM}_\theta-\frac{\xi}{z}+\frac{3e^{-A-\chi}A_\theta}{2})-e^{-A-\chi}}}{z}
    \end{aligned}
\end{equation}

如果用视界边界条件：
\begin{equation}
    \begin{aligned}
        (\frac{f}{z^{3}})^{\prime}&=\frac{\xi_{\theta}}{z^{3}}+\frac{L^{2}}{2z^{4}}e^{-\chi}V+\frac{\mathrm{COM}\textcolor{red}{-\cot\theta(\mathrm{COM}_\theta-\frac{\xi}{z}+\frac{3e^{-A-\chi}A_\theta}{2})-e^{-A-\chi}}}{z^{2}}\\
        zf'-3f&=z\xi_{x}+\frac{L^{2}}{2}e^{-\chi}V+z^2(\mathrm{COM}\textcolor{red}{-\cot\theta(\mathrm{COM}_\theta-\frac{\xi}{z}+\frac{3e^{-A-\chi}A_\theta}{2})-e^{-A-\chi}})\\
        &=2z\xi_{x}\textcolor{red}{+2\cot\theta z\xi}+\frac{L^{2}}{2}e^{-\chi}V+z^2[(\mathrm{COM}-\frac{\xi_x}{z})\textcolor{red}{-\cot\theta([\mathrm{COM}_\theta+\frac{\xi}{z}]+\frac{3e^{-A-\chi}A_\theta}{2})-e^{-A-\chi}}]\\
        &=2z\xi_{x}\textcolor{red}{+2\cot\theta z\xi}+\frac{L^{2}}{2}e^{-\chi}V+z^2[\mathrm{COM}_f\textcolor{red}{-\cot\theta(\mathrm{COM}_{\theta f}+\frac{3e^{-A-\chi}A_\theta}{2})-e^{-A-\chi}}]\\
        \mathrm{COM}_f&=\mathrm{COM}-\frac{\xi_x}{z}\\
        &=\frac{e^{A+\chi}}{4}\xi^{\prime2}-\frac{\xi_{x}^{\prime}}{2}-\frac{e^{-A-\chi}}{4}(2A_x\chi_x+\chi_{x}^2-\phi_x^2)+\frac{(e^{-A-\chi})_{xx}}{2}\\
        &=\frac{e^{A+\chi}}{4}\xi^{\prime2}-\frac{\xi_{x}^{\prime}}{2}+\frac{e^{-A-\chi}}{2}[A_x^2+A_x\chi_x+\frac{1}{2}\chi_{x}^2+\frac{1}{2}\phi_x^2-(A_{xx}+\chi_{xx})]\\
        \textcolor{red}{\mathrm{COM}_{\theta f}}&\textcolor{red}{=\mathrm{COM}_\theta+\frac{\xi}{z}}\\
        &\textcolor{red}{=\frac{\xi'}{2}+\frac{e^{-A-\chi}\chi_\theta}{2}}
    \end{aligned}
\end{equation}

$e^{-A-\chi}E_{xx}-e^{A-\chi}E_{yy}$ gives the evolution equation of $A$:
\begin{equation}
    \begin{aligned}
        z(\frac{\dot{A}}{z})^{\prime}&=\dot{A}^{\prime}-\frac{\dot{A}}{z}\\
        &=\frac{z^{2}}{2}(\frac{fA^{\prime}-\xi A_{\theta}}{z^{2}})^{\prime}+\frac{(e^{-A-\chi}A_{\theta}-\xi A^{\prime})_{\theta}}{2}+\mathrm{COM}\textcolor{red}{+\cot\theta(-\frac{\xi A'}{2}+\frac{\xi'}{2}+\frac{e^{-A-\chi}\chi_\theta}{2}-\frac{\xi}{z})}\\
        &=\frac{z^{2}}{2}(\frac{fA^{\prime}-\xi A_{\theta}}{z^{2}})^{\prime}+\frac{(e^{-A-\chi}A_{\theta}-\xi A^{\prime})_{\theta}}{2}+\mathrm{COM}\textcolor{red}{+\cot\theta(\mathrm{COM}_\theta-\frac{\xi A'}{2})}\\
        &=\frac{(fA^{\prime}-\xi A_{\theta})'}{2}-\frac{fA^{\prime}-\xi A_{\theta}}{z}+\frac{(e^{-A-\chi}A_{\theta}-\xi A^{\prime})_{\theta}}{2}+\mathrm{COM}\textcolor{red}{+\cot\theta(\mathrm{COM}_\theta-\frac{\xi A'}{2})}
    \end{aligned}
\end{equation}

the scalar equation gives the evolution equation of $\phi$:
\begin{equation}
    \begin{aligned}
        z(\frac{\dot{\phi}}{z})^{\prime}&=\dot{\phi}^{\prime}-\frac{\dot{\phi}}{z}\\
        &=\frac{z^{2}}{2}(\frac{f\phi'-\xi\phi_{\theta}}{z^{2}})^{\prime}+\frac{(e^{-A-\chi}\phi_{\theta}-\xi\phi')_{\theta}}{2}-\frac{L^{2}}{2z^{2}}e^{-\chi}V_\phi\textcolor{red}{+\cot\theta(-\frac{\xi \phi'}{2}+\frac{e^{-A-\chi}\phi_\theta}{2})}\\
        &=\frac{(f\phi'-\xi\phi_{\theta})'}{2}-\frac{f\phi'-\xi\phi_{\theta}}{z}+\frac{(e^{-A-\chi}\phi_{\theta}-\xi\phi')_{\theta}}{2}-\frac{L^{2}}{2z^{2}}e^{-\chi}V_\phi\textcolor{red}{+\cot\theta(\frac{e^{-A-\chi}\phi_\theta}{2}-\frac{\xi \phi'}{2})}
    \end{aligned}
\end{equation}

$E_{xn}=E_{xv}-fE_{xz}+\xi E_{xx}$ gives the evolution equation of $\xi$:
\begin{equation}
    \begin{aligned}
        0&=\frac{2}{z}(f_{\theta}+f\chi_{\theta})+(2\xi e^{A+\chi}\xi^{\prime}-f^{\prime}+f[A^{\prime}+\chi^{\prime}])_{\theta}-\xi(e^{A+\chi}\xi^{\prime})_{\theta}-\xi(A+\chi)_{xx}\\
        &\quad+A_{\theta}d_{n}A+\phi_{\theta}d_n\phi-(\dot{A}+\dot{\chi})_{\theta}+(e^{A+\chi}\xi^{\prime})_{v}\\
        &=\frac{2}{z}(f_{\theta}+f\chi_{\theta})+2\xi_{\theta}e^{A+\chi}\xi^{\prime}+(f[A^{\prime}+\chi^{\prime}]-f^{\prime})_{\theta}+\xi(e^{A+\chi}\xi^{\prime})_{\theta}-\xi(A+\chi)_{xx}\\
        &\quad+A_{\theta}d_{n}A+\phi_{\theta}d_n\phi-(\dot{A}+\dot{\chi})_{\theta}+(e^{A+\chi}\xi^{\prime})_{v}\\
        &=\frac{2}{z}(f_{\theta}+f\chi_{\theta})+(2\xi e^{A+\chi}\xi^{\prime}-f^{\prime}+f[A^{\prime}+\chi^{\prime}]-[\dot{A}+\dot{\chi}])_{\theta}-\xi(e^{A+\chi}\xi^{\prime}+A_x+\chi_x)_{\theta}\\
        &\quad+A_{\theta}d_{n}A+\phi_{\theta}d_n\phi+(e^{A+\chi}\xi^{\prime})_{v}\\
        &\quad\textcolor{red}{+\cot\theta(\xi e^{A+\chi}\xi'-3\xi A_\theta+\xi\chi_\theta+2f A'-2\dot{A})-2\xi}\\
        &=\textcolor{red}{+\cot\theta[\xi(e^{A+\chi}\xi'-3A_\theta+\chi_\theta)+2f A'-2\dot{A}]-2\xi}
    \end{aligned}
\end{equation}

where $d_{n}:=\partial_{v}-f\partial_{z}+\xi\partial_{\theta}$.

$E_{vn}=E_{vv}-fE_{vz}+\xi E_{vx}$ gives the evolution equation of $f$:
\begin{equation}
    \begin{aligned}
        0&=(e^{-A-\chi}f_{x}-[f\xi]^{\prime}+\xi^{2}e^{A+\chi}\xi^{\prime}-\xi[d_{n}A+d_{n}\chi]-2d_{n}\xi)_{x}-\xi_x(\dot{A}+\dot{\chi})\\
        &\quad-\frac{2}{z}(\dot{f}+f\dot{\chi})-\dot{A}d_{n}A-\dot{\phi}d_n\phi+\xi(e^{A+\chi}\xi^{\prime})_{v}\\
        &\quad\textcolor{red}{+\cot\theta(\xi^2 e^{A+\chi}\xi'-\xi^2 A_\theta-\xi^2\chi_\theta+f\xi A'+f\xi\chi'-2\xi\dot{\chi}-2\xi\xi_\theta-\xi f'+f\xi'-2\dot{\xi}+e^{-A-\chi}f_\theta)}\\
        &=\textcolor{red}{+\cot\theta(\xi^2[e^{A+\chi}\xi'-(A_\theta+\chi_\theta)]+\xi[f(A'+\chi')-2\dot{\chi}-2\xi_\theta-f']+f\xi'-2\dot{\xi}+e^{-A-\chi}f_\theta)}
    \end{aligned}
\end{equation}

如果把$\theta$的取值范围从$[0,\pi]$扩充到$[-\pi,\pi]$，从方程中可以看出，$\xi$关于$\theta$是奇函数，而$f,\chi,A,\phi$关于$\theta$都是偶函数。另外注意，每求一次导就会让奇偶性交换一次。

\subsubsection{asymptotic behavior}

Assumning $L=1$, we expand the fields near the boundary ($z=0$):
\begin{equation}
    \begin{aligned}
        f&=f_0+f_1 z+f_2 z^2+f_3 z^3+O(z^4)\\
        \chi&=\chi_0+\chi_1 z+\chi_2 z^2+\chi_3 z^3+O(z^4)\\
        \xi&=\xi_0+\xi_1 z+\xi_2 z^2+\xi_3 z^3+O(z^4)\\
        A&=A_0+A_1 z+A_2 z^2+A_3 z^3+O(z^4)\\
        \phi&=\phi_0+\phi_1 z+\phi_2 z^2+\phi_3 z^3+O(z^4)\\
    \end{aligned}
\end{equation}

both $f_i$ are function of $(v,\theta)$, substituting them into equations of motion, we can get
\begin{equation}
    \chi_1=0,\phi_0=0,\xi_1=e^{-A_0-\chi_0}\partial_\theta\chi_0,f_0=e^{-\chi_0},f_1=\textcolor{red}{-\cot\theta\xi_0}-\partial_\theta\xi_0
\end{equation}

after letting $\chi_0=0,\xi_0=0$, we can get
\begin{equation}
    \xi_1=0,f_0=1,f_1=0
\end{equation}
\begin{equation}
    A_1=\partial_v A_0,A_2=0
\end{equation}

after letting $A_0=0$, we can get
\begin{equation}
    A_1=0
\end{equation}
\begin{equation}
    \xi_2=0,f_2=\chi_2\textcolor{red}{+1},\chi_2=\frac{\phi_1^2}{8},\chi_3=\frac{\phi_1\phi_2}{3}
\end{equation}
\begin{equation}
    \phi_3=\partial_v\phi_2-\frac{119\phi_1^3}{360}-\frac{\partial_\theta^2\phi_1}{2}\textcolor{red}{-\cot\theta\frac{\partial_\theta\phi_1}{2}}
\end{equation}
\begin{equation}
    \partial_v\xi_3=\frac{\partial_x f_3}{3}-\partial_x A_3-\frac{2\phi_1\partial_x\phi_2}{9}+\frac{\phi_2\partial_x\phi_1}{9}-\frac{\partial_v\phi_1\partial_x\phi_1}{4}+\frac{\phi_1\partial_{vx}\phi_1}{12}\textcolor{red}{-2\cot\theta A_3}
\end{equation}
\begin{equation}
    \partial_v f_3=\frac{3\partial_x\xi_3}{2}+\frac{\phi_1\partial_v\phi_2}{6}+\frac{2\phi_2\partial_v\phi_1}{3}-\frac{(\partial_v\phi_1)^2}{2}+\frac{(\partial_x\phi_1)^2}{8}+\frac{\phi_1\partial_x^2\phi_1}{8}\textcolor{red}{+\cot\theta(\frac{3\xi_3}{2}+\frac{\phi_1\partial_\theta\phi_1}{8})}
\end{equation}

after letting $\phi_1(v,\theta)=\phi_1$, we can get
\begin{equation}
    \begin{aligned}
        f&=1+(\frac{\phi_1^2}{8}\textcolor{red}{+1})z^2+f_3 z^3+O(z^4)\\
        \chi&=\frac{\phi_1^2}{8}z^2+\frac{\phi_1\phi_2}{3}z^3+O(z^4)\\
        \xi&=\xi_3 z^3+O(z^4)\\
        A&=A_3 z^3+O(z^4)\\
        \phi&=\phi_1 z+\phi_2 z^2+(\partial_v\phi_2-\frac{119\phi_1^3}{360})z^3+O(z^4)\\
    \end{aligned}
\end{equation}
\begin{equation}
    \partial_v\xi_3=\frac{\partial_x f_3}{3}-\partial_x A_3-\frac{2\phi_1\partial_x\phi_2}{9}\textcolor{red}{-2\cot\theta A_3}=\partial_x(\frac{ f_3}{3}-A_3-\frac{2\phi_1\phi_2}{9})\textcolor{red}{-2\cot\theta A_3}
\end{equation}
\begin{equation}
    \partial_v f_3=\frac{3\partial_x\xi_3}{2}+\frac{\phi_1\partial_v\phi_2}{6}\textcolor{red}{+\cot\theta\frac{3\xi_3}{2}}
\end{equation}

\subsection{apparent horizon}

\begin{equation}
    k_0=\xi h_x + \frac{f}{2} + \frac{e^{- A - \chi} h_x^2}{2}
\end{equation}

apparent horizon condition
\begin{equation}
    \begin{aligned}
        0&=\textcolor{red}{\cot\theta\xi}+\xi_x + \xi' h_x - e^{-A-\chi} A_x h_x - e^{-A-\chi} A' h_x^2 - e^{-A-\chi} \chi_x h_x - e^{-A-\chi} \chi' h_x^2\\
        &\quad + e^{-A-\chi} h_{xx} \textcolor{red}{+\cot\theta e^{-A-\chi}h_x} + \frac{f}{h} - \frac{e^{-A-\chi} h_x^2}{h}
    \end{aligned}
\end{equation}

i.e.
\begin{equation}
    (\partial_{x}+h_{x}\partial_{z})(\xi+h_{x}e^{-A-\chi})\textcolor{red}{+\cot\theta(\xi+h_{x}e^{-A-\chi})}=-\frac{1}{h}(f-h_{x}^{2}e^{-A-\chi})
\end{equation}

the evolution of apparent horizon
\begin{equation}
    \begin{aligned}
        &\quad(\partial_{x}+h_{x}\partial_{z})([\partial_{v}+\dot{h}\partial_{z}]\xi+\dot{h}_{x}e^{-A-\chi}+h_{x}[\partial_{v}+\dot{h}\partial_{z}]e^{-A-\chi})+\dot{h}_{x}(\xi+h_{x}e^{-A-\chi})^{\prime}\\
        &\quad\textcolor{red}{+\cot\theta([\partial_{v}+\dot{h}\partial_{z}]\xi+\dot{h}_{x}e^{-A-\chi}+h_{x}[\partial_{v}+\dot{h}\partial_{z}]e^{-A-\chi})}\\
        &=-\frac{1}{h}([\partial_{v}+\dot{h}\partial_{z}]f-2h_{x}\dot{h}_{x}e^{-A-\chi}-h_{x}^{2}[\partial_{v}+\dot{h}\partial_{z}]e^{-A-\chi})+h^{-2}\dot{h}(f-h_{x}^{2}e^{-A-\chi})
    \end{aligned}
\end{equation}

\begin{equation}
    A_0\dot{h}+A_1\dot{h}_x+A_2\dot{h}_{xx}+B=0
\end{equation}

\begin{equation}
    \begin{aligned}
        A_0&=\xi'_{x'}+\left(h_x(e^{-A-\chi})'\right)_{x'}\textcolor{red}{+\cot\theta(\xi+h_{x}e^{-A-\chi})'}+\frac{f'}{h}-\frac{h_x^2}{h}(e^{-A-\chi})'-\frac{f-h_x^2 e^{-A-\chi}}{h^2}\\
        &=(\xi+h_x e^{-A-\chi})'_{x'}\textcolor{red}{+\cot\theta(\xi+h_{x}e^{-A-\chi})'}+\frac{f'}{h}-\frac{h_x^2}{h}(e^{-A-\chi})'-\frac{f-h_x^2 e^{-A-\chi}}{h^2}\\
        A_1&=\xi'+(e^{-A-\chi})_{x'}+h_x(e^{-A-\chi})'+(\xi+h_x e^{-A-\chi})'\textcolor{red}{+\cot\theta e^{-A-\chi}}-\frac{2h_x}{h}e^{-A-\chi}\\
        &=(e^{-A-\chi})_{x'}+2(\xi+h_x e^{-A-\chi})'\textcolor{red}{+\cot\theta e^{-A-\chi}}-\frac{2h_x}{h}e^{-A-\chi}\\
        A_2&=e^{-A-\chi}\\
        B&=(\dot{\xi}+h_x\partial_v(e^{-A-\chi}))_{x'}\textcolor{red}{+\cot\theta(\dot{\xi}+h_x\partial_v(e^{-A-\chi}))}+\frac{\dot{f}-h_x^2\partial_v(e^{-A-\chi})}{h}
    \end{aligned}
\end{equation}

\subsubsection{check}
the metric
\begin{equation}
    ds^2=\frac{L^2}{z^2}(-[f(v,z,\theta)e^{-\chi(v,z,\theta)}-e^{A(v,z,\theta)}\xi(v,z,\theta)^2]dv^2-2e^{-\chi}dvdz-2\xi e^Advd\theta+e^Ad\theta^2+e^{-A}\sin^2\theta d\varphi^2)
\end{equation}

the inverse of the metric
\begin{equation}
    \left[\begin{matrix}
        0 & - z^{2} e^{\chi} & 0 & 0\\- z^{2} e^{\chi} & z^{2} f e^{\chi} & - z^{2} \xi e^{\chi} & 0\\0 & - z^{2} \xi e^{\chi} & z^{2} e^{- A} & 0\\0 & 0 & 0 & \frac{z^{2} e^{A}}{\sin^{2}{\left(\theta \right)}}
    \end{matrix}\right]
\end{equation}

\begin{equation}
    (dz)_\mu=(0,1,0,0),\quad(dz)^\mu=g^{\mu\nu}(dz)_\nu=-z^2 e^{\chi}\cdot(1,-f,\xi,0)\propto(\partial_n)^\mu
\end{equation}

\begin{equation}
    E^{zz},{E^z}_b,E_{ab}
\end{equation}

\begin{equation}
    {E^z}_\varphi=0,E_{v\varphi}=0,E_{\theta\varphi}=0
\end{equation}

$\partial_v E_{zz}$ gives the evolution equation of $\chi$:
\begin{equation}
    \dot{\chi}'=\frac{z}{2}(A'\dot{A}'+\phi'\dot{\phi}')
\end{equation}

$e^{-A-\chi}E_{\theta\theta}+e^{A-\chi}E_{\varphi\varphi}/\sin^2\theta$ gives rise to
\begin{equation}
    \begin{aligned}
        0&=- \xi \frac{\partial}{\partial \theta} A \frac{\partial}{\partial z} A - \xi \frac{\partial}{\partial \theta} \phi \frac{\partial}{\partial z} \phi + 2 \xi \frac{\partial^{2}}{\partial z\partial \theta} \chi + \frac{\xi \cos{\left(\theta \right)} \frac{\partial}{\partial z} A}{\sin{\left(\theta \right)}} + \frac{f \left(\frac{\partial}{\partial z} A\right)^{2}}{2} - f \frac{\partial^{2}}{\partial z^{2}} \chi + \frac{f \left(\frac{\partial}{\partial z} \phi\right)^{2}}{2}\\
        &\quad- e^{A} e^{\chi} \left(\frac{\partial}{\partial z} \xi\right)^{2} - \frac{\partial}{\partial v} A \frac{\partial}{\partial z} A - \frac{\partial}{\partial z} A \frac{\partial}{\partial \theta} \xi + \frac{\partial}{\partial \theta} \chi \frac{\partial}{\partial z} \xi - \frac{\partial}{\partial z} \chi \frac{\partial}{\partial z} f - \frac{\partial}{\partial v} \phi \frac{\partial}{\partial z} \phi + \frac{\partial^{2}}{\partial z^{2}} f\\
        &\quad+ \textcolor{red}{2 \frac{\partial^{2}}{\partial z\partial v} \chi} - \frac{\partial^{2}}{\partial z\partial \theta} \xi - \frac{\cos{\left(\theta \right)} \frac{\partial}{\partial z} \xi}{\sin{\left(\theta \right)}} + e^{- A} e^{- \chi} \frac{\partial}{\partial \theta} A \frac{\partial}{\partial \theta} \chi + e^{- A} e^{- \chi} \left(\frac{\partial}{\partial \theta} \chi\right)^{2} - e^{- A} e^{- \chi} \frac{\partial^{2}}{\partial \theta^{2}} \chi\\
        &\quad- \frac{e^{- A} e^{- \chi} \cos{\left(\theta \right)} \frac{\partial}{\partial \theta} \chi}{\sin{\left(\theta \right)}} + \frac{2 \xi \cos{\left(\theta \right)}}{z \sin{\left(\theta \right)}} + \frac{2 \frac{\partial}{\partial \theta} \xi}{z} - \frac{4 \frac{\partial}{\partial z} f}{z} + \frac{V{\left(\phi \right)} e^{- \chi}}{z^{2}} + \frac{6 f}{z^{2}}
    \end{aligned}
\end{equation}

\end{document}